\chapter{锁、条件变量与信号量}

锁的作用是保护共享状态，让临界区在同一时间只被一个线程执行。锁不是性能敌人，也不是正确性的全部。正确使用锁，需要明确保护对象、锁顺序、临界区长度和等待条件。

\section{从共享状态到互斥锁}

每把锁都应该有清晰的保护范围。一个队列锁保护 head、tail、size 和 buffer 状态；一个计数器锁保护 value；一个对象锁保护对象不变量。锁如果没有对应的不变量，就很容易出现“看起来加锁了但仍然错”的情况。

临界区应尽量短，但不能为了短而破坏不变量。把状态检查放在锁外，可能引入竞态；把耗时 I/O 放在锁内，会扩大阻塞范围。工程上要在正确性和粒度之间找到边界。

\begin{lstlisting}[language=C++,caption={锁保护队列不变量}]
class QueueState {
public:
    void push(std::int32_t value) {
        std::lock_guard<std::mutex> lock(this->mutex_);
        this->values_.push_back(value);
    }

private:
    std::mutex mutex_;
    std::vector<std::int32_t> values_;
};
\end{lstlisting}

这段代码的关键不是 \code{push_back} 这一行，而是锁和不变量的关系：\code{values_} 的内部结构不能被多个线程同时修改。如果未来再加入 \code{size_}、\code{head_} 或统计字段，它们也必须在同一把锁保护下保持一致。

写锁代码时，应先写“锁保护的不变量”，再写代码。比如有界队列的不变量可以是：\code{0 <= size <= capacity}，\code{head} 和 \code{tail} 总是落在环形数组范围内，\code{size == 0} 表示不能 pop，\code{size == capacity} 表示不能 push。互斥锁保护的是这些关系，不只是保护某一个变量。若 \code{size} 在锁内改，\code{head} 在锁外改，这个不变量就被拆碎了。

\topic{锁粒度}

粗粒度锁容易证明正确，但并发度低。细粒度锁提高并发度，却增加锁顺序、死锁和状态拆分复杂度。分片锁是一种常见折中：按 key、队列分段或线程组拆成多把锁，让无关操作并行，热点操作仍受保护。

锁粒度应从访问模式出发。如果所有请求都访问同一个热点 key，分片没有帮助；如果请求天然分散，分片锁可以显著减少竞争。锁优化前要先测量等待时间和竞争比例。

\section{从等待谓词到条件变量}

条件变量用于等待某个状态变真。它必须和互斥锁配合，并在循环里检查谓词。原因是线程可能被虚假唤醒，也可能在被唤醒后发现条件已经被其他线程消费。正确模式是“持锁检查条件，不满足则等待，醒来继续检查”。

\begin{lstlisting}[language=C++,caption={条件变量等待谓词}]
std::unique_lock<std::mutex> lock(this->mutex_);
this->not_empty_.wait(lock, [this]() {
    return !this->values_.empty();
});
const std::int32_t value = this->values_.front();
\end{lstlisting}

这里谓词属于受锁保护的状态。通知本身不携带数据，数据在共享状态里。若不用循环检查谓词，就可能在虚假唤醒或竞争消费后读取空队列。

条件变量最重要的句子是：等待的是谓词，不是通知。通知可能在等待前发生，也可能唤醒多个线程，也可能出现虚假唤醒。只要谓词在锁保护下检查，程序就正确；只依赖“我收到过一次 notify”，程序就脆弱。生产者消费者队列里，\code{not_empty} 的谓词是队列非空，\code{not_full} 的谓词是队列未满。两个谓词对应同一组受锁保护状态。

\topic{信号量和资源计数}

信号量表示可用资源数量。它适合限制并发度、实现 bounded buffer 或控制连接池。信号量和互斥锁不同：互斥锁保护状态的一致性，信号量表达资源额度。混用时要小心锁顺序，避免死锁。

\section{锁顺序、异常和性能路径}

死锁通常需要互斥、持有并等待、不可抢占和循环等待四个条件。破坏循环等待的常见方法是规定全局锁顺序。复杂系统里，锁顺序文档和锁级别检查很重要，因为死锁往往在低概率路径中出现。

\topic{优先级反转和 convoy}

优先级反转发生在低优先级线程持有锁，高优先级线程等待它，而中等优先级线程持续占用 CPU，导致高优先级线程无法前进。锁 convoy 则是多个线程排队等待同一锁，锁释放后唤醒和抢占成本让系统吞吐下降。这些问题说明锁不仅是用户态对象，也和调度器行为相互作用。

工程系统要避免在锁内做不可控工作，例如磁盘 I/O、网络调用、复杂回调和用户插件。锁内执行时间越不可预测，越容易出现尾延迟问题。

\topic{锁的性能路径}

mutex 通常有三条路径。第一条是无竞争 fast path，用户态原子操作成功，成本较低。第二条是短暂竞争路径，线程可能自旋几轮，等待锁持有者很快释放。第三条是阻塞路径，线程进入内核等待，涉及 futex、调度和唤醒。性能分析要判断当前主要走哪条路径。一个低竞争 mutex 没必要改成复杂无锁结构；一个长期进入阻塞路径的热点锁，则需要拆分状态、缩短临界区或改变算法。

临界区长度不只由代码行数决定。锁内分配内存、写日志、调用用户回调、访问远端 NUMA 内存、发生 page fault、触发异常路径，都会让持锁时间变长。锁优化报告应列出锁内可能的慢操作，而不只是说“这个锁很粗”。

\topic{读写锁和 RCU 直觉}

读多写少场景中，读写锁允许多个读者并发，但写者需要独占。它适合读临界区不太长、写入不太频繁的结构。若读者很多且读路径必须极轻，RCU 的思想更进一步：读者在不阻塞写者的情况下读取旧版本，写者发布新版本后，等所有旧读者离开再回收旧对象。

RCU 的难点不在“读很快”这个口号，而在对象生命周期。旧对象什么时候没人再读，什么时候可以释放，如何处理写者之间的同步，都是必须证明的问题。后面讲无锁内存回收时，会看到相同主题。

\section{有界队列、关闭和背压}

Compute Systems Engine 的第一个同步结构应是有界队列，因为它把互斥、条件变量、背压和关闭语义放在同一个小系统里。有界队列不是一个简单容器，它表达了生产者和消费者之间的容量契约：队列未满时生产者可以放入，队列非空时消费者可以取出，队列满时生产者必须等待或失败，队列关闭后等待者必须能醒来并退出。

先写不变量，再写代码。一个环形有界队列至少有这些不变量：容量大于零；\code{head} 和 \code{tail} 总在 \code{[0, capacity)} 范围内；\code{size} 总在 \code{[0, capacity]} 范围内；\code{size == 0} 时没有元素可取；\code{size == capacity} 时没有空槽可写；\code{head} 指向下一个可读槽；\code{tail} 指向下一个可写槽。互斥锁保护的是这组关系，而不是某一个整数。

只要写代码时先写不变量，很多错误会提前暴露。例如 \code{push} 修改 \code{tail} 但忘了增加 \code{size}，\code{try_pop} 读取元素后忘了通知 \code{not_full}，\code{size} 用无符号类型导致下溢，析构时仍有线程阻塞，关闭后生产者继续写入，这些都不是语法问题，而是不变量和状态转换没有定义清楚。

\topic{从 try 到 blocking 的接口分层}

队列接口不应只有一种操作。常见接口至少分成三类：非阻塞尝试、阻塞等待和关闭。非阻塞 \code{try_push} 或 \code{try_pop} 立即返回，适合事件循环或自定义调度；阻塞 \code{push} 或 \code{pop} 等待谓词变真，适合普通生产者消费者；关闭 \code{close} 唤醒等待者，告诉它们系统不再接受新任务或不再产生新元素。

如果队列没有关闭语义，线程池很难优雅退出。worker 可能永远阻塞在空队列上，主线程析构线程池时只能强行终止进程或依赖外部哨兵值。哨兵值在泛型队列中也不可靠，因为每种元素类型未必都有安全哨兵。正确做法是把关闭作为队列状态的一部分，受同一把锁保护，并在关闭时通知所有等待者。

\begin{lstlisting}[language=C++,caption={有关闭语义的有界队列接口草图}]
class BoundedQueue {
public:
    explicit BoundedQueue(std::int32_t capacity);

    bool try_push(std::int32_t value);
    bool push(std::int32_t value);
    std::optional<std::int32_t> pop();
    void close();

private:
    std::vector<std::int32_t> buffer_;
    std::int32_t head_ = 0;
    std::int32_t tail_ = 0;
    std::int32_t size_ = 0;
    bool closed_ = false;
    std::mutex mutex_;
    std::condition_variable not_empty_;
    std::condition_variable not_full_;
};
\end{lstlisting}

这里 \code{push} 返回 bool，是为了表达关闭后不能再写入；\code{pop} 返回 optional，是为了表达关闭且队列为空时消费者应退出。这个接口比单纯抛异常更适合作为线程池内部队列，因为关闭不是异常，而是正常生命周期事件。当然，业务层接口也可以选择抛异常，但内部状态机仍要明确。

\topic{等待谓词的完整写法}

生产者等待的谓词不是“收到 not full 通知”，而是“队列未满或已经关闭”。消费者等待的谓词不是“收到 not empty 通知”，而是“队列非空或已经关闭”。关闭必须出现在谓词里，否则关闭时即使 \code{notify_all}，线程醒来后发现队列仍满或仍空，又会继续睡回去，线程池无法退出。

\begin{lstlisting}[language=C++,caption={push 的状态转换}]
bool BoundedQueue::push(std::int32_t value) {
    std::unique_lock<std::mutex> lock(this->mutex_);
    this->not_full_.wait(lock, [this]() {
        return this->closed_ ||
               this->size_ < static_cast<std::int32_t>(this->buffer_.size());
    });
    if (this->closed_) {
        return false;
    }

    this->buffer_[static_cast<std::size_t>(this->tail_)] = value;
    this->tail_ = (this->tail_ + 1) %
                  static_cast<std::int32_t>(this->buffer_.size());
    ++this->size_;
    this->not_empty_.notify_one();
    return true;
}
\end{lstlisting}

\begin{lstlisting}[language=C++,caption={pop 的状态转换}]
std::optional<std::int32_t> BoundedQueue::pop() {
    std::unique_lock<std::mutex> lock(this->mutex_);
    this->not_empty_.wait(lock, [this]() {
        return this->closed_ || this->size_ > 0;
    });
    if (this->size_ == 0) {
        return std::nullopt;
    }

    const std::int32_t value =
        this->buffer_[static_cast<std::size_t>(this->head_)];
    this->head_ = (this->head_ + 1) %
                  static_cast<std::int32_t>(this->buffer_.size());
    --this->size_;
    this->not_full_.notify_one();
    return value;
}
\end{lstlisting}

这两个片段体现了条件变量的关键原则：状态改变和谓词检查都在同一把锁保护下；通知只用于让等待者重新检查谓词；关闭是状态机的一部分。注意通知可以在持锁时调用，也可以在释放锁后调用，不同写法有不同唤醒和竞争细节。教材阶段更重要的是先保证谓词正确，再讨论微小性能差异。

\topic{锁内代码的边界}

有界队列的锁内代码应只做状态转换，不应执行用户回调、复杂分配、日志格式化或长时间 I/O。若 \code{push} 在锁内调用用户提供的处理函数，队列锁就保护了不可控代码，任何慢操作都会阻塞所有生产者和消费者。若 \code{pop} 在锁内执行任务，线程池全局队列会退化成串行执行器。

这条原则可以推广到所有锁。锁内只维护不变量，锁外做耗时工作。若必须在锁内移动复杂对象，应考虑预分配、移动语义、异常安全和对象析构成本。C++ 中对象析构也可能执行代码，如果最后一个引用在锁内释放，析构函数可能做不可控工作。高质量并发代码会尽量让重析构发生在锁外。

\topic{异常安全和状态回滚}

C++ 队列如果存储复杂对象，\code{push} 可能在移动或拷贝元素时抛异常。若异常发生在 \code{tail} 和 \code{size} 已经修改之后，不变量就会损坏。教学示例用 \code{std::int32_t} 可以避开这个问题，但正式教材必须说明边界：泛型队列需要把可能抛异常的操作和状态提交点分开，或者要求元素移动不抛异常，或者使用预构造槽位和 placement new  carefully 管理生命周期。

状态机里要有提交点。对于 \code{push}，元素成功放入槽位并且 \code{tail/size} 更新后，操作才对消费者可见；对于 \code{pop}，元素被取出并且 \code{head/size} 更新后，槽位才重新可用。若中间步骤可能失败，就要保证失败前队列仍保持旧状态，或失败后能回滚。锁只能提供互斥，不能自动提供异常安全。

\topic{信号量实现 bounded buffer 的视角}

有界队列也可以用两个信号量表达：一个信号量记录空槽数量，一个信号量记录已有元素数量，再用互斥保护环形索引。生产者先获取空槽，再写入；消费者先获取元素，再读取。这个模型把“资源数量”和“状态互斥”分开，适合帮助读者理解信号量。

但信号量版本同样需要关闭语义。若消费者阻塞在元素信号量上，关闭时如何唤醒？若生产者阻塞在空槽信号量上，关闭时如何退出？标准信号量本身不携带“关闭”状态，需要额外状态和唤醒策略。条件变量版本把谓词写在代码里，更适合教学；信号量版本更接近资源额度控制。两者没有绝对优劣，关键是状态机是否完整。

\topic{锁顺序和组合操作}

单个队列容易证明，多个锁组合就会出现锁顺序问题。假设线程池有全局队列锁、worker 本地队列锁和指标锁。如果 worker 持有本地队列锁时尝试拿全局队列锁，而另一个线程持有全局队列锁时尝试拿本地队列锁，就可能死锁。工程系统应定义锁级别：例如先全局配置锁，再队列锁，再指标锁；或者避免同时持有多把锁，用消息和局部复制拆开组合操作。

组合操作还会诱导“为了省事加一把大锁”。大锁让正确性简单，却可能把整个运行时串行化。合理做法是先用粗锁实现 reference，写清不变量和测试；再根据测量拆分锁，而不是一开始就写复杂细锁。拆锁时，每拆一把都要重新写不变量：哪些状态仍由旧锁保护，哪些状态由新锁保护，跨锁关系如何保持。

\topic{futex 慢路径和用户态 fast path}

Linux 上常见 mutex 和 condition variable 都会尽量把无竞争路径留在用户态。无竞争加锁可能只是一次原子状态改变；竞争时才进入 futex 让线程睡眠。条件变量等待也会释放 mutex 并进入等待，唤醒后重新竞争 mutex。这个设计说明锁不是一上来就系统调用，系统调用通常是竞争或等待的慢路径。

阅读 futex 相关内容时，要抓住三个点。第一，用户态变量保存同步状态，内核只在需要睡眠和唤醒时参与。第二，等待必须和用户态状态检查配合，避免丢失唤醒。第三，唤醒不等于立即执行，被唤醒线程还要进入调度器、被安排到 CPU，并重新竞争锁。因此高竞争条件变量的成本不只是 \code{notify_one}，还包括调度和再次竞争。

\topic{互斥锁保护的是不变量}

很多人把 mutex 理解成“保护变量”，这句话不够精确。mutex 保护的是一组变量之间的不变量。以环形队列为例，\code{head}、\code{tail}、\code{size} 和 \code{buffer} 不是四个孤立变量，而是共同满足容量边界、满空状态、元素位置和可用槽位这些关系。只把 \code{size} 改成原子并不能让队列线程安全，因为 \code{size} 和 \code{head/tail} 的关系仍然会被并发修改破坏。

写并发代码时，第一步应写不变量表。哪些字段属于同一个状态机，哪些字段只用于指标，哪些字段可以近似，哪些字段必须精确。锁的粒度也由不变量决定。如果两个字段永远一起变化，拆成两把锁可能反而制造跨锁一致性问题；如果两个字段从语义上独立，一把大锁可能只是偷懒。锁设计不是“粗锁慢、细锁快”的简单比较，而是不变量边界和竞争形状的平衡。

\begin{lstlisting}[numbers=none]
bounded queue invariant:
  0 <= size <= capacity
  head and tail are valid positions
  size == 0 means no readable slot
  size == capacity means no writable slot
  readable slots form one circular interval from head
  writable slots are the complement interval
\end{lstlisting}

有了这张表，代码审查就更具体。每个会修改状态的函数，都必须在持锁状态下从一个合法状态转换到另一个合法状态。每个等待谓词，都必须对应一个状态条件，而不是对应一次通知。每个关闭路径，都必须说明关闭后哪些操作仍允许，哪些操作返回失败，哪些等待者必须被唤醒。

\topic{条件变量不是事件}

条件变量最常见的误解，是把它当成事件通知系统。正确理解是：条件变量让线程在某个谓词不满足时睡眠，并在可能满足时醒来重新检查。通知本身不保存状态；如果通知发生时没有等待者，通知不会留在条件变量里等后来者。状态必须保存在受 mutex 保护的变量中。

这条原则可以解释为什么等待必须写成循环或谓词版本。线程可能被虚假唤醒；也可能被唤醒后别的线程先抢到锁并消费了资源；也可能关闭状态改变但队列仍空。醒来不代表条件成立，只代表“应该重新检查”。把 \code{if} 写成 \code{while} 不是形式主义，而是条件变量语义的一部分。

\begin{lstlisting}[language=C++,caption={条件变量等待的是谓词，不是通知次数}]
std::optional<std::int32_t> pop_wait() {
    std::unique_lock<std::mutex> lock(this->mutex_);
    this->not_empty_.wait(lock, [this]() {
        return this->closed_ || this->size_ > 0;
    });

    if (this->size_ == 0) {
        return std::nullopt;
    }

    const std::int32_t value =
        this->buffer_[static_cast<std::size_t>(this->head_)];
    this->head_ = (this->head_ + 1) %
                  static_cast<std::int32_t>(this->buffer_.size());
    --this->size_;
    this->not_full_.notify_one();
    return value;
}
\end{lstlisting}

这段代码里，\code{closed_} 和 \code{size_} 都在同一把锁下读写。若关闭函数只设置 \code{closed_} 却不 \code{notify_all}，等待线程可能永远睡眠；若关闭函数通知但不把 \code{closed_} 放进谓词，线程醒来后又睡回去。状态和通知必须配套。

\topic{关闭语义是并发容器的一等接口}

很多教学队列只写 \code{push} 和 \code{pop}，但线程池真正需要的是关闭语义。关闭不是析构的附属动作，而是容器状态机的一部分。一个队列关闭后，生产者是否还能写入？消费者是否能继续取出已经存在的元素？空队列关闭后消费者返回什么？满队列关闭后生产者如何退出？这些问题决定 worker 能否优雅停止。

常见策略是：关闭后拒绝新 \code{push}；已经在队列中的元素仍允许被 \code{pop} 取出；当队列为空且关闭时，\code{pop} 返回空值；关闭函数唤醒所有生产者和消费者。这种策略适合线程池 drain：不再接收新任务，但已提交任务可以执行完。另一种策略是 cancel：关闭后直接丢弃队列中剩余任务。它适合进程退出或错误恢复，但必须明确通知调用者哪些任务未执行。没有声明的关闭策略，会让上层业务在退出、异常和超时时出现悬挂。

\begin{lstlisting}[numbers=none]
queue close policy:
  open + not full: push succeeds
  open + full: producer waits
  closed: push fails
  closed + nonempty: pop drains existing item
  closed + empty: pop returns null
\end{lstlisting}

线程池的 \code{shutdown}、\code{shutdown_now}、\code{wait} 和析构函数应围绕这套语义设计。析构函数里悄悄阻塞等待所有任务完成，可能让用户在持锁析构时死锁；析构函数直接丢任务，又可能破坏业务语义。教材应鼓励读者把关闭写成显式 API，而不是把复杂语义藏在析构里。

\topic{锁粒度的演化}

一个复杂系统最稳妥的路线，通常是先写粗锁 reference。粗锁版本把不变量集中在一处，容易测试，容易与串行 reference 对比。等到测量证明锁确实是瓶颈，再拆成细锁、分片锁、读写锁或无锁结构。直接从细锁开始，很容易得到一个既不快也不正确的系统。

拆锁时要写迁移计划。第一步，确认锁保护的全部状态和不变量。第二步，把状态按独立性分组，设计新锁的保护范围。第三步，找出跨组操作，例如移动任务、关闭队列、汇总指标。第四步，为跨组操作定义锁顺序或避免同时持锁。第五步，保留粗锁版本作为 reference 或测试 oracle。拆锁不是把一个 mutex 机械替换成多个 mutex，而是重新设计状态机边界。

读写锁也不是默认更快。读多写少且读临界区较长时，它可能有收益；读临界区很短时，读写锁本身的元数据和公平策略可能比普通 mutex 更重；写者频繁时，读写锁还可能造成写者饥饿或读者抖动。若读路径只需要不可变快照，RCU 或 copy-on-write 可能比读写锁更适合。同步原语必须和访问模式匹配。

\topic{锁竞争的观测}

锁竞争不能只靠感觉。应用指标可以记录等待次数、等待总时间、最长等待、持锁时间、队列长度和唤醒次数。系统工具可以观察 futex 调用、调度切换、CPU 利用率和 off-CPU 时间。若一个程序 CPU 使用率很低但延迟很高，可能大量时间在等待锁或 I/O；若 CPU 使用率很高但吞吐不升，可能在自旋或 CAS 失败。

持锁时间和等待时间要分开。持锁时间长，说明临界区里工作太多；等待时间长，可能因为竞争高，也可能因为持锁线程被调度走。若临界区很短但等待时间长，要看线程数是否远超核心、是否有优先级反转、是否发生 NUMA 迁移。锁问题是执行图问题，不是单行代码问题。

在 Linux 上，可以结合 \cmd{perf}、\cmd{strace -e futex}、调度 trace 和应用内部统计。教材不要求每个读者都能开启所有权限，但要让读者知道证据链：用户态等待指标说明哪里堵，futex 说明是否进入内核睡眠，调度指标说明线程是否被抢占，CPU 拓扑说明 cache line 迁移是否昂贵。

\topic{优先级反转和公平性}

锁还涉及调度语义。优先级反转指高优先级线程等待低优先级线程持有的锁，而低优先级线程又迟迟得不到 CPU。实时系统中，这会造成严重延迟。某些系统通过优先级继承缓解：持锁的低优先级线程临时继承等待者的高优先级，尽快释放锁。普通服务端虽然不一定使用实时优先级，也会遇到类似现象：关键请求等待后台任务持有的锁，尾延迟被放大。

公平锁按等待顺序唤醒线程，可以减少饥饿，但可能降低吞吐，因为它限制了“刚好在 CPU 上的线程再次获得锁”的机会。非公平锁吞吐可能更高，但某些线程可能长期等待。条件变量唤醒一个还是唤醒全部，也有公平和惊群权衡。高性能系统要根据业务目标选择：批处理更关心总吞吐，交互服务更关心尾延迟和饥饿。

\topic{RAII 和异常边界}

C++ 中锁必须用 RAII 管理。手写 \code{lock} 和 \code{unlock} 容易在异常、早返回和多分支中漏解锁。\code{std::lock_guard} 适合简单临界区，\code{std::unique_lock} 适合条件变量、延迟加锁和手动解锁。教材示例应让读者形成习惯：锁对象的作用域就是临界区，作用域越小越好，锁外做耗时工作。

RAII 也不能替你设计不变量。一个函数里拿了锁，但把引用返回给外部，让外部在锁释放后修改受保护状态，仍然错误。一个类内部用 mutex，但复制构造没有定义好，也可能把锁和状态复制出奇怪语义。并发类通常应禁用复制，或者显式定义移动和共享语义。锁只是同步工具，不是封装边界的替代品。

\topic{生产者消费者的完整测试}

生产者消费者测试不能只检查几个值是否按顺序出来。完整测试至少包含容量为一、容量为多、单生产者单消费者、多生产者多消费者、生产者快消费者慢、消费者快生产者慢、关闭空队列、关闭非空队列、满队列关闭、重复关闭、析构前关闭、压力运行和 ThreadSanitizer。每个测试对应一个状态机边界。

还要测试没有丢失和没有重复。可以让每个生产者生成带 producer id 和 sequence 的任务，消费者记录所有任务，最终检查总数、唯一性和每个生产者内部顺序。如果队列声明全局 FIFO，就检查全局顺序；如果只声明每生产者顺序或无序，就不要测试不存在的保证。接口保证必须和测试一致。

\begin{lstlisting}[numbers=none]
queue stress test data:
  producer id: identifies writer
  sequence: monotonically increasing per producer
  payload: optional checksum
  consumer log: records every consumed item
  final check: no loss, no duplicate, declared order holds
\end{lstlisting}

性能测试也要有边界。线程创建成本不应混进队列操作；输入生成不应混进消费；日志输出不应放在热循环；关闭时间和正常处理时间要分开。否则队列性能报告很容易测到无关成本。

\topic{读写锁的适用边界}

读写锁允许多个读者并发进入，写者独占进入。它看起来很适合“读多写少”，但真实收益取决于读临界区长度、写入频率、实现公平性和缓存行为。若读临界区很短，读写锁内部维护读者计数和状态的成本可能超过普通 mutex；若写者频繁，读者并发优势会被写者阻塞抵消；若实现偏向读者，写者可能饥饿；若实现偏向写者，读者尾延迟可能变差。

读写锁特别适合读操作较长、读者之间不修改共享状态、写操作低频且能接受等待的结构，例如较大的配置表、路由快照或只读索引。它不适合每次请求都要读写的小状态，也不适合读路径本身会更新统计字段的场景。很多“读操作”其实包含隐藏写，例如更新 LRU、增加引用计数、记录访问时间，这会破坏读者并发。

若读多写少结构可以使用不可变快照，RCU 或 copy-on-write 可能比读写锁更好。读者读取一个稳定指针，不加锁；写者复制一份新结构，修改后发布；旧结构延迟回收。这样读路径更轻，但写路径复制成本和回收复杂度更高。读写锁、RCU、copy-on-write 都是同一问题的不同取舍：读者成本、写者成本、一致性和回收。

\topic{锁分解案例：统计服务}

假设一个统计服务维护三类状态：全局配置、每个 worker 的计数、最近错误列表。朴素实现可以用一把大锁保护所有状态，正确性最简单，但每次更新计数都和读取配置、写错误列表互相阻塞。拆锁前，应先写状态关系。配置很少更新，读多；计数高频写，读少；错误列表低频写，偶尔查询。它们的访问模式不同，适合不同同步策略。

一种改进是配置使用快照或读写锁，计数使用每 worker 分片，错误列表使用单独 mutex。这样高频计数不再抢配置锁，错误查询也不阻塞计数热路径。拆锁后要检查跨状态操作。例如“更新配置并清空统计”是否需要同时触碰配置和计数？如果需要，就要定义锁顺序或改成两阶段操作。拆锁的难点不在把 mutex 数量变多，而在找出跨状态不变量。

\begin{lstlisting}[language=C++,caption={按访问模式拆分状态的结构形态}]
struct alignas(64) WorkerCounter {
    std::uint64_t accepted = 0;
    std::uint64_t rejected = 0;
};

class RuntimeStats {
public:
    explicit RuntimeStats(std::int32_t worker_count)
        : counters_(static_cast<std::size_t>(worker_count)) {}

    void record_accept(std::int32_t worker_id) {
        this->counters_[static_cast<std::size_t>(worker_id)].accepted += 1;
    }

private:
    std::vector<WorkerCounter> counters_;
    mutable std::mutex error_mutex_;
    std::vector<std::string> recent_errors_;
};
\end{lstlisting}

这段代码仍然省略了并发读取总数的同步语义。若只有 worker 写、实验结束后汇总，可以不加锁；若运行中监控线程读取，就要定义近似还是精确。如果要求精确实时快照，需要额外同步，分片计数的写入优势会被读取语义部分抵消。同步设计总是从读写语义出发。

\topic{锁顺序表}

多个锁共存后，项目需要锁顺序表。锁顺序表不是形式主义，它让代码审查能发现潜在死锁。例如规定：配置锁先于任务图锁，任务图锁先于队列锁，队列锁先于指标锁。任何代码如果反向获取，就必须重构或证明不会并发。没有锁顺序表，死锁常常在很久后由一次看似无关的功能触发。

锁顺序表还要说明哪些锁不能在持有时调用用户回调、不能做 I/O、不能等待 future。最危险的不是持有两把锁，而是持锁进入未知代码。用户回调可能再进入运行时，拿另一把锁，或者阻塞等待当前锁保护的状态。高质量库通常在锁内只修改内部状态，复制出需要处理的数据，然后锁外调用回调。

\topic{条件变量和多谓词}

真实队列常有多个等待条件：不空、不满、关闭、取消、优先级任务到达、队列水位下降。把所有条件混成一个布尔变量会让状态机混乱。更好的方式是用明确状态字段和谓词函数。例如生产者等待“有空间或关闭”，消费者等待“有任务或关闭”，drain 等待“未完成计数为零”。每个谓词都应能用当前状态直接判断。

多个条件变量可以减少无关唤醒，但也增加复杂度。一个 bounded queue 常用 \code{not_empty} 和 \code{not_full} 两个条件变量；线程池可能还有 \code{idle_cv} 用于 wait；关闭时通常要 \code{notify_all} 所有相关条件变量。通知遗漏会导致悬挂，通知过多会导致惊群。先保证状态机正确，再用指标判断是否需要优化唤醒策略。

\topic{Semaphore 和资源额度}

信号量适合表达资源额度。数据库连接池、并发请求上限、I/O in-flight 限额、令牌桶都可以用信号量思想理解。它和 mutex 的区别是：mutex 保护临界区不变量，semaphore 控制有多少个许可。把二者混淆会导致设计不清。一个线程拿到许可后，仍可能需要 mutex 修改共享队列。

信号量还要处理取消和关闭。线程等待许可时，如果上层请求超时，等待能否被取消？系统关闭时，等待者如何醒来？许可已经拿到但任务失败，是否归还？这些都是资源状态机的一部分。信号量不是“比条件变量简单”的替代品，它只是把数量资源表达得更直接。

\topic{锁保护对象的文档格式}

并发类应在代码附近写清每个锁保护哪些字段，以及哪些字段是原子或线程本地。注释不需要冗长，但要能帮助维护者判断修改是否安全。比如：“\code{mutex_} protects head, tail, size, buffer, closed”，“\code{stats_} is worker-owned and read after join”，“\code{closed_} is only read while holding mutex”。这种注释比泛泛写“thread-safe”有用得多。

接口文档也要写清线程安全级别。某个对象是否允许多个线程同时调用 \code{push}？是否允许一个线程析构时其他线程仍在调用？\code{close} 是否幂等？\code{wait} 是否可以和 \code{submit} 并发？这些问题如果不写，调用者会按自己的理解使用，最终把内部不变量打破。

\topic{案例：关闭线程池}

关闭线程池是锁和条件变量的综合案例。线程池有任务队列、未完成计数、worker 列表和状态。调用 \code{shutdown} 后，\code{submit} 应拒绝新任务，worker 应继续取出已接受任务，队列为空且状态为 shutting down 时退出。\code{wait_idle} 应等待未完成计数为零，但不一定关闭线程池。析构应确保 worker 已 join。

错误实现常见于三个地方。第一，只设置 \code{closed} 但没有 \code{notify_all}，睡眠 worker 不醒。第二，队列空就让 worker 退出，但 \code{submit} 仍可能并发加入任务。第三，任务抛异常后未完成计数没有减少，\code{wait} 永久阻塞。正确实现需要一张状态机，而不是几处布尔判断。

\topic{案例：有界队列和内存预算}

有界队列的容量可以按任务数，也可以按字节。线程池任务通常按数量限制；I/O batch 队列更需要按字节限制。若每个任务大小差别很大，只限制数量可能导致内存爆炸。锁保护的状态也会更复杂：不仅有 \code{size}，还有 \code{bytes_in_queue}。生产者等待谓词变成“任务数未满且字节数未超，或关闭”。

这个案例说明锁保护的是不变量集合。任务数、字节数、head/tail、closed 共同决定队列状态。拆开任何一个，都可能破坏背压语义。

\topic{锁章验收清单}

本章完成后，项目里的每个锁都应能回答：保护哪些字段，临界区是否调用用户代码，是否可能同时持有多把锁，锁顺序是什么，等待谓词是什么，关闭时谁被唤醒，异常路径是否保持不变量，测试是否覆盖容量一、关闭、并发压力和重复关闭。若回答不出来，这把锁就是维护风险。

\topic{锁章节总结}

锁的价值是让复杂不变量变得可证明，条件变量的价值是让线程在谓词不满足时睡眠，信号量的价值是表达资源额度。它们不是低级替代品，而是并发系统的基础工具。先用锁写正确 reference，再用测量决定是否拆锁、分片、读写锁、RCU 或无锁。把锁视为失败，只会诱导更难证明的错误代码。

\topic{同步原语决策表}

选择同步原语时，先看状态形状。一个复杂不变量，用 mutex；等待某个谓词，用 condition variable；控制资源额度，用 semaphore；读多写少且读临界区较长，用读写锁或快照；只需要统计，用 relaxed atomic；单生产者单消费者队列，可以考虑 SPSC ring；多生产者多消费者且关闭语义复杂，先用 mutex bounded queue。选择同步原语不是按“高级程度”，而是按状态语义。

\topic{案例：锁内析构}

C++ 中最后一个 \code{std::shared_ptr} 在锁内离开作用域，可能触发对象析构。析构函数如果写日志、释放大内存、关闭文件或调用回调，就把不可控工作放进锁内。高质量代码会在锁内把对象移动到局部变量，释放锁后再让局部变量析构。这个细节常常影响尾延迟。

锁内析构说明临界区不只包含显式语句，也包含对象生命周期副作用。审查锁内代码时，要看构造、析构、分配、回调和异常，不只看表面几行。

\topic{案例：等待谓词丢失}

一个常见 bug 是消费者等待 \code{not_empty}，但关闭时只设置 \code{closed} 并通知 \code{not_full}。结果消费者如果正睡在 \code{not_empty}，就永远不醒。另一个 bug 是等待谓词只写 \code{size} 大于零，不包含 \code{closed}；关闭后消费者醒来，发现 \code{size} 仍为零，又睡回去。正确代码必须让每个等待者的谓词覆盖所有能让它退出等待的状态。

这个案例看似简单，却是线程池停不下来的常见根源。并发代码的退出路径和成功路径一样重要。


\topic{从一次具体事故开始}

有界队列在关闭时偶发卡死，另一个计数表版本虽然加了锁，却在线程数增加后几乎不扩展。这不是抽象练习，而是后续所有机制的入口。锁 API 本身没错，错在没有写清锁保护的不变量、条件变量等待的谓词和多锁组合的顺序。如果这一层只写成术语列表，读者会知道很多名词，却不知道面对这类现象时先固定什么、再观察什么、最后改什么。

本章把运行问题固定为有界队列、共享计数表和线程池状态。这件事的价值是让所有推导都落在同一条数据流上：输入如何进入系统，状态在哪里改变，输出何时对外可见，失败发生后哪些东西还能相信。主失败可以压缩成一句话：只保护代码片段而不保护不变量，会造成丢唤醒、死锁、关闭悬挂和扩展性下降。后面的每个机制都要回到这个失败，而不是各讲各的。

\begin{lstlisting}[numbers=none]
mutex protects: buffer, head, tail, size, closed
not_empty waits: size > 0 or closed
not_full waits: size < capacity or closed
\end{lstlisting}

阅读本章时要先画这条链路，而不是先背概念。链路中每个箭头都表示一次边界跨越：从文件到内存，从线程到队列，从局部状态到共享状态，从临时结果到提交结果。边界越多，隐含假设越多；隐含假设越多，优化时越容易把正确性、性能和恢复搅在一起。

本章的目标不是给出一个万能框架，而是训练一种判断顺序：先写 reference，再暴露朴素版本，再沿着失败路径引入机制，最后用实验和审查表确定结论边界。如果一个优化不能说明它改变了哪条边界，就先不要把它合进贯穿项目。

\topic{朴素方案为什么会被写出来}

第一版哪个函数会访问共享变量就在哪个函数加锁；关闭时设置 closed 标志，但没有唤醒所有等待者。这个写法并不是一开始就荒谬。它符合小程序直觉：对象少、状态近、调用栈清楚、失败罕见，所有问题都能在一个调试会话里看见。经典教材写到这里不能直接嘲笑朴素方案，因为读者需要先理解它在哪些条件下成立。

朴素方案的真正问题，是它把多个边界藏了起来。控制状态和业务数据混在一起，临时结果和提交结果混在一起，测量边界和语义边界混在一起，线程资源和任务数量混在一起。代码短只是表面现象；代价是没有地方记录不变量，也没有地方插入观测点。

把这个方案放大到有界队列、共享计数表和线程池状态，第一个变化是输入规模让偶发路径变成常态，第二个变化是并发让顺序假设失效，第三个变化是系统资源让等待和数据移动变成主成本。所以改进不是在原函数里继续加分支，而是把边界显式化。

教材里的朴素版本必须保留。它一方面作为 correctness oracle 的对照，另一方面作为解释机制的反面样本。没有朴素版本，最终设计就像凭空出现；没有反例，最终设计又会被误解成永远正确。

\topic{失败路径一：结果错误不是性能问题}

消费者在空队列上等待，生产者在满队列上等待，close 只改标志不通知，线程可能永远睡眠。这类问题通常不会稳定复现。小输入没有触发，单线程没有触发，甚至同一批输入换一个调度顺序才触发。因此第一反应不能是调线程数，也不能是继续优化热循环，而是先把语义锚点拿出来。

排查时按四步走。第一步固定输入和随机种子；第二步运行串行 reference，保存可比较摘要；第三步让优化版本输出同样的摘要；第四步定位第一个不同的边界。摘要不能只写总数，还要包含错误分类、分片身份、attempt 身份、输出 checksum 和阶段状态。

\begin{lstlisting}[numbers=none]
t0 run reference and save checksum
t1 run naive implementation on the same input
t2 compare record count, error count, key count, output checksum
t3 stop at the first boundary that diverges
\end{lstlisting}

这里的关键是“第一个不同的边界”。如果 reference 与优化版本在解析阶段已经不同，就不要继续讨论 cache、SIMD、线程池或分布式提交。系统分析要避免把下游症状当成根因。

\topic{失败路径二：吞吐提升可能掩盖资源失控}

分片转移时两个队列按调用顺序上锁，不同线程形成等待环；全局 map 锁还把所有 key 更新串行化。这类失败常常伴随一个迷惑性信号：平均吞吐看起来变好，尾延迟、内存峰值、恢复时间或输出可信度却变坏。高质量教材必须把这类反例写进正文，否则读者会误以为所有优化只要更快就是更好。

资源失控要按队列、内存、文件、线程和网络五类检查。队列看水位和等待，内存看峰值和分配次数，文件看临时产物和持久化点，线程看 runnable 与 blocked，网络或消息看 in-flight 与重试。单个总耗时无法解释这些状态。

\begin{lstlisting}[numbers=none]
observe:
  queue_depth, producer_wait_ns, consumer_wait_ns
  resident_bytes, allocations, temporary_files
  runnable_threads, blocked_threads, retry_count
  committed_outputs, stale_outputs, checksum_errors
\end{lstlisting}

如果某个指标没有采集，就在报告里明确写成未验证风险。不要把没有观测到的路径写成不存在，也不要用当前机器的一次稳定运行去替代边界证明。

\topic{把机制落成状态和接口}

本章要求每把锁都能回答：保护哪个不变量，等待哪个谓词，谁负责通知，异常时状态如何恢复。状态和接口是机制进入工程的地方。概念如果不能落成字段、状态、函数边界、错误码或实验指标，就很难在代码审查里被检查。

\begin{lstlisting}[language=C++,caption={BoundedQueueInvariant 的最小教学结构}]
template <typename T>
struct BoundedQueueState {
    std::vector<T> buffer;
    std::uint64_t head = 0;
    std::uint64_t tail = 0;
    std::uint64_t size = 0;
    bool closed = false;
};
\end{lstlisting}

这段结构不是要替代真实项目代码，而是把本章最容易被忽略的边界写出来。字段名应回答一个具体问题：身份是什么，状态属于谁，哪次 attempt 有效，哪些字节已经被处理，哪个 checksum 能证明输出没有被误读。

接口设计还有一个原则：控制面字段不要被热路径随手修改，数据面批量结构不要夹带恢复协议。前者让状态机可审查，后者让硬件看到规则的数据流。若两者混在一起，代码会同时难以优化和难以恢复。

\topic{机制推导：不变量}

先看一个具体问题：锁到底保护变量还是保护关系。在有界队列、共享计数表和线程池状态中，这个问题会沿着输入、执行、同步、I/O 或提交边界继续传播，不会停留在一个局部函数里。

朴素写法通常是：哪个变量会被多线程访问就给哪个变量加锁。它吸引人的地方是短、直观、容易在小样本上通过测试；危险也在这里，因为小样本会替它隐藏资源、调度、数据分布和失败路径。

放大以后会出现的信号是：多个变量之间的关系可能在锁外被破坏。这个信号未必是崩溃，也可能是吞吐平台期、p99 抖动、结果偶尔变化、队列持续变长、重启后状态不一致，或者某个分片长期拖尾。

此时引入不变量，不是为了增加术语，而是为了建立一个可检查模型：锁保护的是不变量，例如 head、tail、size 和 buffer 内容的一致关系。模型必须同时解释正确性和成本。只解释为什么快，不解释为什么可信，不能进入贯穿项目；只解释为什么正确，却完全无视缓存、调度、I/O 或网络，也不符合第二册目标。

把模型写成不变量时，要问四个问题：哪些状态只能由一个拥有者修改，哪些结果允许重试，哪些数据可以批量搬运，哪些错误必须外显。围绕不变量的代码审查，就从这四个问题开始，而不是从实现看起来是否高级开始。

观察方式应围绕假设设计：为每个临界区写前置条件和后置条件。实验要固定输入和环境，保留 reference，一次只改变一个变量。如果某项工具在当前手机或 Linux 环境没有权限，就写出预期命令、缺失事件和结论限制。

最后要写边界：锁范围太大降低并发，太小破坏语义。边界不是附录里的保守声明，而是正文的一部分。读者应该知道什么时候使用它，什么时候退回简单方案，什么时候需要换一个尺度重新测。

本节验收可以用一句话判断：能否从朴素版本推导到改进版本，并能说清新增字段、新增状态或新增实验到底保护了什么。如果只能给最终代码，却解释不了这条推导链，说明机制还没有真正学会。

\topic{机制推导：临界区粒度}

先看一个具体问题：一个大锁简单，为什么有时不够好。在有界队列、共享计数表和线程池状态中，这个问题会沿着输入、执行、同步、I/O 或提交边界继续传播，不会停留在一个局部函数里。

朴素写法通常是：用一个 mutex 保护整个对象。它吸引人的地方是短、直观、容易在小样本上通过测试；危险也在这里，因为小样本会替它隐藏资源、调度、数据分布和失败路径。

放大以后会出现的信号是：低竞争时清楚，高竞争时所有操作串行。这个信号未必是崩溃，也可能是吞吐平台期、p99 抖动、结果偶尔变化、队列持续变长、重启后状态不一致，或者某个分片长期拖尾。

此时引入临界区粒度，不是为了增加术语，而是为了建立一个可检查模型：粒度选择是在正确性、复杂度和争用之间取舍。模型必须同时解释正确性和成本。只解释为什么快，不解释为什么可信，不能进入贯穿项目；只解释为什么正确，却完全无视缓存、调度、I/O 或网络，也不符合第二册目标。

把模型写成不变量时，要问四个问题：哪些状态只能由一个拥有者修改，哪些结果允许重试，哪些数据可以批量搬运，哪些错误必须外显。围绕临界区粒度的代码审查，就从这四个问题开始，而不是从实现看起来是否高级开始。

观察方式应围绕假设设计：比较全局锁、分片锁和读写锁。实验要固定输入和环境，保留 reference，一次只改变一个变量。如果某项工具在当前手机或 Linux 环境没有权限，就写出预期命令、缺失事件和结论限制。

最后要写边界：分锁必须定义锁顺序，否则死锁风险上升。边界不是附录里的保守声明，而是正文的一部分。读者应该知道什么时候使用它，什么时候退回简单方案，什么时候需要换一个尺度重新测。

本节验收可以用一句话判断：能否从朴素版本推导到改进版本，并能说清新增字段、新增状态或新增实验到底保护了什么。如果只能给最终代码，却解释不了这条推导链，说明机制还没有真正学会。

\topic{机制推导：条件变量谓词}

先看一个具体问题：为什么 wait 必须放在 while 谓词里。在有界队列、共享计数表和线程池状态中，这个问题会沿着输入、执行、同步、I/O 或提交边界继续传播，不会停留在一个局部函数里。

朴素写法通常是：收到 notify 就认为条件成立。它吸引人的地方是短、直观、容易在小样本上通过测试；危险也在这里，因为小样本会替它隐藏资源、调度、数据分布和失败路径。

放大以后会出现的信号是：虚假唤醒、竞争唤醒和关闭状态会让线程醒来后条件仍不成立。这个信号未必是崩溃，也可能是吞吐平台期、p99 抖动、结果偶尔变化、队列持续变长、重启后状态不一致，或者某个分片长期拖尾。

此时引入条件变量谓词，不是为了增加术语，而是为了建立一个可检查模型：条件变量等待的是状态谓词，不是通知本身。模型必须同时解释正确性和成本。只解释为什么快，不解释为什么可信，不能进入贯穿项目；只解释为什么正确，却完全无视缓存、调度、I/O 或网络，也不符合第二册目标。

把模型写成不变量时，要问四个问题：哪些状态只能由一个拥有者修改，哪些结果允许重试，哪些数据可以批量搬运，哪些错误必须外显。围绕条件变量谓词的代码审查，就从这四个问题开始，而不是从实现看起来是否高级开始。

观察方式应围绕假设设计：构造多消费者和 close 场景验证。实验要固定输入和环境，保留 reference，一次只改变一个变量。如果某项工具在当前手机或 Linux 环境没有权限，就写出预期命令、缺失事件和结论限制。

最后要写边界：notify 不是消息队列，不能携带业务事件。边界不是附录里的保守声明，而是正文的一部分。读者应该知道什么时候使用它，什么时候退回简单方案，什么时候需要换一个尺度重新测。

本节验收可以用一句话判断：能否从朴素版本推导到改进版本，并能说清新增字段、新增状态或新增实验到底保护了什么。如果只能给最终代码，却解释不了这条推导链，说明机制还没有真正学会。

\topic{机制推导：丢唤醒}

先看一个具体问题：先 notify 后 wait 会不会丢事件。在有界队列、共享计数表和线程池状态中，这个问题会沿着输入、执行、同步、I/O 或提交边界继续传播，不会停留在一个局部函数里。

朴素写法通常是：把通知当成存储的信号。它吸引人的地方是短、直观、容易在小样本上通过测试；危险也在这里，因为小样本会替它隐藏资源、调度、数据分布和失败路径。

放大以后会出现的信号是：若状态没有记录，等待者可能永远睡眠。这个信号未必是崩溃，也可能是吞吐平台期、p99 抖动、结果偶尔变化、队列持续变长、重启后状态不一致，或者某个分片长期拖尾。

此时引入丢唤醒，不是为了增加术语，而是为了建立一个可检查模型：正确模式是先修改受锁保护状态，再 notify，让等待者检查状态。模型必须同时解释正确性和成本。只解释为什么快，不解释为什么可信，不能进入贯穿项目；只解释为什么正确，却完全无视缓存、调度、I/O 或网络，也不符合第二册目标。

把模型写成不变量时，要问四个问题：哪些状态只能由一个拥有者修改，哪些结果允许重试，哪些数据可以批量搬运，哪些错误必须外显。围绕丢唤醒的代码审查，就从这四个问题开始，而不是从实现看起来是否高级开始。

观察方式应围绕假设设计：用压力测试和超时测试暴露悬挂。实验要固定输入和环境，保留 reference，一次只改变一个变量。如果某项工具在当前手机或 Linux 环境没有权限，就写出预期命令、缺失事件和结论限制。

最后要写边界：靠 sleep 修复丢唤醒是错误做法。边界不是附录里的保守声明，而是正文的一部分。读者应该知道什么时候使用它，什么时候退回简单方案，什么时候需要换一个尺度重新测。

本节验收可以用一句话判断：能否从朴素版本推导到改进版本，并能说清新增字段、新增状态或新增实验到底保护了什么。如果只能给最终代码，却解释不了这条推导链，说明机制还没有真正学会。

\topic{工程边界：信号量}

先看一个具体问题：资源槽位为什么不一定需要完整队列锁。在有界队列、共享计数表和线程池状态中，这个问题会沿着输入、执行、同步、I/O 或提交边界继续传播，不会停留在一个局部函数里。

朴素写法通常是：用 mutex 加计数器手写等待。它吸引人的地方是短、直观、容易在小样本上通过测试；危险也在这里，因为小样本会替它隐藏资源、调度、数据分布和失败路径。

放大以后会出现的信号是：容易遗漏释放路径和异常安全。这个信号未必是崩溃，也可能是吞吐平台期、p99 抖动、结果偶尔变化、队列持续变长、重启后状态不一致，或者某个分片长期拖尾。

此时引入信号量，不是为了增加术语，而是为了建立一个可检查模型：信号量表示可用资源数量，适合连接数、缓冲槽和并发限流。模型必须同时解释正确性和成本。只解释为什么快，不解释为什么可信，不能进入贯穿项目；只解释为什么正确，却完全无视缓存、调度、I/O 或网络，也不符合第二册目标。

把模型写成不变量时，要问四个问题：哪些状态只能由一个拥有者修改，哪些结果允许重试，哪些数据可以批量搬运，哪些错误必须外显。围绕信号量的代码审查，就从这四个问题开始，而不是从实现看起来是否高级开始。

观察方式应围绕假设设计：测试 acquire、release、超时和异常路径。实验要固定输入和环境，保留 reference，一次只改变一个变量。如果某项工具在当前手机或 Linux 环境没有权限，就写出预期命令、缺失事件和结论限制。

最后要写边界：信号量不表达队列内容所有权，仍需数据结构同步。边界不是附录里的保守声明，而是正文的一部分。读者应该知道什么时候使用它，什么时候退回简单方案，什么时候需要换一个尺度重新测。

本节验收可以用一句话判断：能否从朴素版本推导到改进版本，并能说清新增字段、新增状态或新增实验到底保护了什么。如果只能给最终代码，却解释不了这条推导链，说明机制还没有真正学会。

\topic{工程边界：死锁}

先看一个具体问题：两个锁都正确为什么组合后会卡死。在有界队列、共享计数表和线程池状态中，这个问题会沿着输入、执行、同步、I/O 或提交边界继续传播，不会停留在一个局部函数里。

朴素写法通常是：每个函数按自己方便的顺序加锁。它吸引人的地方是短、直观、容易在小样本上通过测试；危险也在这里，因为小样本会替它隐藏资源、调度、数据分布和失败路径。

放大以后会出现的信号是：线程之间形成等待环。这个信号未必是崩溃，也可能是吞吐平台期、p99 抖动、结果偶尔变化、队列持续变长、重启后状态不一致，或者某个分片长期拖尾。

此时引入死锁，不是为了增加术语，而是为了建立一个可检查模型：死锁来自资源获取顺序和不可抢占等待。模型必须同时解释正确性和成本。只解释为什么快，不解释为什么可信，不能进入贯穿项目；只解释为什么正确，却完全无视缓存、调度、I/O 或网络，也不符合第二册目标。

把模型写成不变量时，要问四个问题：哪些状态只能由一个拥有者修改，哪些结果允许重试，哪些数据可以批量搬运，哪些错误必须外显。围绕死锁的代码审查，就从这四个问题开始，而不是从实现看起来是否高级开始。

观察方式应围绕假设设计：画等待图，定义全局锁顺序。实验要固定输入和环境，保留 reference，一次只改变一个变量。如果某项工具在当前手机或 Linux 环境没有权限，就写出预期命令、缺失事件和结论限制。

最后要写边界：try lock 和超时只能缓解，不能替代设计。边界不是附录里的保守声明，而是正文的一部分。读者应该知道什么时候使用它，什么时候退回简单方案，什么时候需要换一个尺度重新测。

本节验收可以用一句话判断：能否从朴素版本推导到改进版本，并能说清新增字段、新增状态或新增实验到底保护了什么。如果只能给最终代码，却解释不了这条推导链，说明机制还没有真正学会。

\topic{工程边界：异常安全}

先看一个具体问题：临界区里抛异常会不会留下半更新状态。在有界队列、共享计数表和线程池状态中，这个问题会沿着输入、执行、同步、I/O 或提交边界继续传播，不会停留在一个局部函数里。

朴素写法通常是：手动 lock 后在多个分支 unlock。它吸引人的地方是短、直观、容易在小样本上通过测试；危险也在这里，因为小样本会替它隐藏资源、调度、数据分布和失败路径。

放大以后会出现的信号是：异常或早返回导致锁不释放或状态半更新。这个信号未必是崩溃，也可能是吞吐平台期、p99 抖动、结果偶尔变化、队列持续变长、重启后状态不一致，或者某个分片长期拖尾。

此时引入异常安全，不是为了增加术语，而是为了建立一个可检查模型：RAII 和先构造后提交能保持异常路径清楚。模型必须同时解释正确性和成本。只解释为什么快，不解释为什么可信，不能进入贯穿项目；只解释为什么正确，却完全无视缓存、调度、I/O 或网络，也不符合第二册目标。

把模型写成不变量时，要问四个问题：哪些状态只能由一个拥有者修改，哪些结果允许重试，哪些数据可以批量搬运，哪些错误必须外显。围绕异常安全的代码审查，就从这四个问题开始，而不是从实现看起来是否高级开始。

观察方式应围绕假设设计：写抛异常的单元测试验证状态不变量。实验要固定输入和环境，保留 reference，一次只改变一个变量。如果某项工具在当前手机或 Linux 环境没有权限，就写出预期命令、缺失事件和结论限制。

最后要写边界：异常安全不是只靠 lock guard，状态更新顺序也要设计。边界不是附录里的保守声明，而是正文的一部分。读者应该知道什么时候使用它，什么时候退回简单方案，什么时候需要换一个尺度重新测。

本节验收可以用一句话判断：能否从朴素版本推导到改进版本，并能说清新增字段、新增状态或新增实验到底保护了什么。如果只能给最终代码，却解释不了这条推导链，说明机制还没有真正学会。

\topic{工程边界：锁性能路径}

先看一个具体问题：锁慢是因为内核调用还是 cache line 争用。在有界队列、共享计数表和线程池状态中，这个问题会沿着输入、执行、同步、I/O 或提交边界继续传播，不会停留在一个局部函数里。

朴素写法通常是：看到 mutex 就认为一定进内核。它吸引人的地方是短、直观、容易在小样本上通过测试；危险也在这里，因为小样本会替它隐藏资源、调度、数据分布和失败路径。

放大以后会出现的信号是：低竞争 fast path 可能很快，高竞争才进入等待和唤醒。这个信号未必是崩溃，也可能是吞吐平台期、p99 抖动、结果偶尔变化、队列持续变长、重启后状态不一致，或者某个分片长期拖尾。

此时引入锁性能路径，不是为了增加术语，而是为了建立一个可检查模型：锁性能由竞争、临界区长度、唤醒策略和 cache line 所有权共同决定。模型必须同时解释正确性和成本。只解释为什么快，不解释为什么可信，不能进入贯穿项目；只解释为什么正确，却完全无视缓存、调度、I/O 或网络，也不符合第二册目标。

把模型写成不变量时，要问四个问题：哪些状态只能由一个拥有者修改，哪些结果允许重试，哪些数据可以批量搬运，哪些错误必须外显。围绕锁性能路径的代码审查，就从这四个问题开始，而不是从实现看起来是否高级开始。

观察方式应围绕假设设计：记录等待次数、持锁时间和上下文切换。实验要固定输入和环境，保留 reference，一次只改变一个变量。如果某项工具在当前手机或 Linux 环境没有权限，就写出预期命令、缺失事件和结论限制。

最后要写边界：不要为了避免 mutex 直接写错误 atomic 协议。边界不是附录里的保守声明，而是正文的一部分。读者应该知道什么时候使用它，什么时候退回简单方案，什么时候需要换一个尺度重新测。

本节验收可以用一句话判断：能否从朴素版本推导到改进版本，并能说清新增字段、新增状态或新增实验到底保护了什么。如果只能给最终代码，却解释不了这条推导链，说明机制还没有真正学会。

\topic{案例推导：先有问题，再有答案}

本章案例不直接给最终实现。每个案例都按同一条教学链路展开：先描述输入和失败，再写朴素版本为什么自然，接着指出第一处证据，最后才引入改进。这样读者能学会推导，而不是只记住答案。

\topic{案例：有界队列 close}

场景是：生产者和消费者都可能在关闭时等待。先把它缩到可以手工检查的规模，再扩到能触发真实性能或恢复问题的规模。小规模负责证明语义，大规模负责暴露成本，两者缺一不可。
朴素版本是：只设置 closed 标志。它必须保留在仓库里，名字可以直接标成 baseline 或 naive。保留它不是为了上线，而是为了让反例可复现。
改进版本是：在锁内更新状态并 notify all，让 push 和 pop 都检查谓词。改进说明要写清楚它改变了哪条路径：少搬了哪些字节，少争了哪条 cache line，少等了哪个队列，减少了哪些重复请求，或者把哪个半提交状态变成了可恢复状态。
验收重点是：验证等待中关闭、空队列关闭、满队列关闭。报告里至少有一组反例。没有反例的案例很容易变成宣传文字，读者也无法判断机制的边界。

\topic{案例：分片锁 map}

场景是：多个 key 更新同一张计数表。先把它缩到可以手工检查的规模，再扩到能触发真实性能或恢复问题的规模。小规模负责证明语义，大规模负责暴露成本，两者缺一不可。
朴素版本是：整张表一个 mutex。它必须保留在仓库里，名字可以直接标成 baseline 或 naive。保留它不是为了上线，而是为了让反例可复现。
改进版本是：按 shard 加锁并定义合并阶段。改进说明要写清楚它改变了哪条路径：少搬了哪些字节，少争了哪条 cache line，少等了哪个队列，减少了哪些重复请求，或者把哪个半提交状态变成了可恢复状态。
验收重点是：比较热点 key 和均匀 key 下的争用。报告里至少有一组反例。没有反例的案例很容易变成宣传文字，读者也无法判断机制的边界。

\topic{案例：双锁转移}

场景是：把任务从一个队列移动到另一个队列。先把它缩到可以手工检查的规模，再扩到能触发真实性能或恢复问题的规模。小规模负责证明语义，大规模负责暴露成本，两者缺一不可。
朴素版本是：先锁源再锁目标或按调用者顺序锁。它必须保留在仓库里，名字可以直接标成 baseline 或 naive。保留它不是为了上线，而是为了让反例可复现。
改进版本是：按地址或 id 定义全局锁顺序。改进说明要写清楚它改变了哪条路径：少搬了哪些字节，少争了哪条 cache line，少等了哪个队列，减少了哪些重复请求，或者把哪个半提交状态变成了可恢复状态。
验收重点是：压力测试不应悬挂。报告里至少有一组反例。没有反例的案例很容易变成宣传文字，读者也无法判断机制的边界。

\topic{实验矩阵：让结论可复现}

容量设为一，同时启动多个生产者和消费者；在空、满、半满、等待中分别调用 close。实验要把正常输入、边界输入、压力输入和错误输入分开。正常输入证明功能还在，边界输入证明不变量，压力输入暴露成本，错误输入验证恢复。

围绕不变量、临界区粒度、条件变量谓词、丢唤醒，不要把所有变量一次改完。比如改变布局时固定线程数，改变线程数时固定输入分布，改变批大小时固定写入设备状态，改变重试策略时固定故障注入脚本。变量克制是性能报告可信的前提。

\begin{lstlisting}[numbers=none]
experiment matrix:
  input: normal / boundary / pressure / faulty
  version: reference / naive / improved / counterexample
  metrics: correctness / throughput / latency / resource / recovery
  evidence: command / environment / raw sample / conclusion limit
\end{lstlisting}

报告中的数字要能追溯到命令。只贴一张结果表不够；还要写 CPU 或设备环境、编译选项、输入生成方式、是否预热、是否绑定线程、是否清理 page cache、是否启用故障注入。做不到的项目要写成未验证风险。

如果改进版本更慢，也不是失败。负结果常常比正结果更能训练判断力。它可能说明瓶颈不在这里，输入规模太小，机制成本超过收益，或者朴素方案在当前边界下已经足够。工程能力包括知道什么时候不改。

\topic{Linux 源码阅读：从用户态现象倒推对象}

本章 Linux 阅读入口保持窄范围：\filepath{kernel/futex/core.c}、\filepath{kernel/locking/mutex.c}、\filepath{kernel/locking/semaphore.c}。阅读目的不是把内核子系统背下来，而是把一个用户态现象映射到内核对象和状态变化上。

阅读笔记建议固定四列：用户态操作、内核对象、状态变化、可观察指标。用户态操作可以是等待、缺页、写文件、调度、计时或消息收发；内核对象可以是 task、VMA、page、inode、wait queue、bio、socket buffer 或计数器。

源码里的数据结构要比函数名更重要。链表、红黑树、位图、引用计数、状态枚举、等待队列、页表和缓存结构，决定了机制如何落地。读者不需要一次读完整个内核，但要能解释一个最小现象的因果链。

如果当前 Termux 或设备内核无法直接对应源码路径，也要写清版本边界。源码阅读提供机制参考，运行时证据仍来自当前机器上的实验。把二者混为一谈，会让报告看起来权威，实际不可复现。

\topic{排查演练：把事故走完一遍}

排查从具体差异开始：遇到卡死先画等待图，标出每个线程持有什么锁、等待什么谓词；不要用 sleep 或超时掩盖丢唤醒。这一步要慢，不要跳。越早跳到熟悉的工具，越容易把系统带到错误方向。

排查记录应保留时间线。第一行写事故输入和版本，第二行写 reference 摘要，第三行写朴素版本摘要，第四行写第一次分歧，后面才写指标和修复。时间线能防止复盘变成事后故事。

\begin{lstlisting}[numbers=none]
debug timeline:
  1. freeze input and version
  2. run reference and save digest
  3. run naive or optimized version
  4. compare first divergent boundary
  5. inspect metrics for that boundary
  6. patch invariant and rerun matrix
\end{lstlisting}

修复时只改一个边界。若同时改数据布局、线程数、队列容量、批大小和提交协议，即使结果变好也无法解释为什么。高质量教材的案例要能被读者复做，而不是只让读者相信作者经验。

\topic{项目验收边界}

贯穿项目不是把每章代码堆到一起，而是让每章新增一个可证明的能力。能力必须有 reference、朴素实现、改进实现、实验脚本、指标输出和失败注入中的至少几项。

\begin{lstlisting}[numbers=none]
project checkpoints:
  1. 为 bounded queue 写不变量注释
  2. 实现 close 语义和 notify all
  3. 加入锁等待指标
  4. 为分片计数表比较全局锁和分片锁
  5. 写死锁等待图练习
\end{lstlisting}

每个检查点都要进入仓库。只在正文里说“应该实现”不够；读者需要看到代码边界、测试边界和报告边界。若某个检查点受当前设备限制，提交降级说明，而不是把它从质量标准里删除。

项目报告最后写三段：本章新增能力解决了什么失败；新增能力引入了什么成本；哪些结论能迁移，哪些结论必须在新机器、新输入或新故障条件下重测。这个复盘会把整本书连成一条主线。

\topic{深度补充：完整走读、反例和实验手册}

现在把本章再走深一层。前面的正文已经给出事故入口：有界队列在关闭时偶发卡死，另一个计数表版本虽然加了锁，却在线程数增加后几乎不扩展。。这一节不再添加新的名词，而是把一次真实排查拆成可重复的学习动作。读者要能拿着这一节，在自己的机器上复现一个小版本，再把结论迁移到贯穿项目。

第一步仍然是固定任务：有界队列、共享计数表和线程池状态。固定任务不是为了限制想象，而是为了让所有机制共享同一组输入、同一条状态链、同一个 reference 和同一套指标。如果每讲一个概念就换一个例子，读者很难看到概念之间的因果关系；如果一直围绕同一条链路推进，读者能看到一个失败怎样把下一个机制推出来。

第二步是拆出最小可复现实验。最小实验不等于玩具实验，它必须保留会触发本章核心失败的形状：有足够大的输入，有明确的边界条件，有可比较的 reference，有一个朴素版本，有一个改进版本，还有一个能证明边界的反例。缺少任意一项，实验都会变成看起来能跑、实际无法解释的演示。

\begin{lstlisting}[numbers=none]
minimum reproducible study:
  1. freeze input and reference digest
  2. run baseline with stage metrics
  3. inject one pressure or failure condition
  4. change one mechanism only
  5. compare correctness, cost, and recovery
  6. write the counterexample before the conclusion
\end{lstlisting}

第三步是把结果写成证据链，而不是写成心得。证据链的形状很固定：现象是什么，怀疑哪条边界，为什么朴素版本会在这里失败，改进版本保护了哪个不变量，哪些指标支持这个判断，哪些情况没有覆盖。这比直接写“使用某技术后性能提升”更慢，但它能防止误导，也能让后续章节复用结论。

\topic{单点分析：不变量}

围绕不变量，不要先记定义，先问它解决哪一个具体卡点：锁到底保护变量还是保护关系。在本章的有界队列、共享计数表和线程池状态里，这个卡点通常不会独立出现，而是和输入规模、数据分布、执行顺序、资源上限或故障注入一起出现。

朴素版本是：哪个变量会被多线程访问就给哪个变量加锁。它在小输入下往往能通过，因为小输入没有把隐藏假设压穿。但是教材不能停在“它错了”这句话上，要说明它为什么一开始会被写出来：它减少了状态，靠调用栈维持顺序，靠当前机器的资源余量掩盖等待，靠人工观察代替指标。

一旦规模或失败形状变化，就会出现：多个变量之间的关系可能在锁外被破坏。排查这类信号时，先不要改实现。先把它翻译成不变量问题：是不是同一个状态被多个拥有者写，是否某个输出位置没有唯一归属，是否某个对象生命周期跨过了异步边界，是否某个提交点缺少身份和校验。

机制模型是：锁保护的是不变量，例如 head、tail、size 和 buffer 内容的一致关系。把模型落到代码时，至少要能指出一个字段、一个状态转换、一个测试和一个指标。字段让语义进入结构体，状态转换让路径可审查，测试让反例固定，指标让性能判断不靠猜。

实验观察是：为每个临界区写前置条件和后置条件。这里要避免一个常见错误：看到工具输出后直接写结论。正确做法是先写假设，再写为什么这个指标能支持或否定假设。若指标和假设没有关系，就不要把它放进报告。

边界是：锁范围太大降低并发，太小破坏语义。边界要写在正文里，因为工程设计通常不是“会不会”，而是“在什么条件下值得”。一个机制可能在大输入、高并发、失败恢复中必要，却在单线程小工具中完全多余。能说清这个取舍，才说明真正理解了机制。

\topic{单点分析：临界区粒度}

围绕临界区粒度，不要先记定义，先问它解决哪一个具体卡点：一个大锁简单，为什么有时不够好。在本章的有界队列、共享计数表和线程池状态里，这个卡点通常不会独立出现，而是和输入规模、数据分布、执行顺序、资源上限或故障注入一起出现。

朴素版本是：用一个 mutex 保护整个对象。它在小输入下往往能通过，因为小输入没有把隐藏假设压穿。但是教材不能停在“它错了”这句话上，要说明它为什么一开始会被写出来：它减少了状态，靠调用栈维持顺序，靠当前机器的资源余量掩盖等待，靠人工观察代替指标。

一旦规模或失败形状变化，就会出现：低竞争时清楚，高竞争时所有操作串行。排查这类信号时，先不要改实现。先把它翻译成不变量问题：是不是同一个状态被多个拥有者写，是否某个输出位置没有唯一归属，是否某个对象生命周期跨过了异步边界，是否某个提交点缺少身份和校验。

机制模型是：粒度选择是在正确性、复杂度和争用之间取舍。把模型落到代码时，至少要能指出一个字段、一个状态转换、一个测试和一个指标。字段让语义进入结构体，状态转换让路径可审查，测试让反例固定，指标让性能判断不靠猜。

实验观察是：比较全局锁、分片锁和读写锁。这里要避免一个常见错误：看到工具输出后直接写结论。正确做法是先写假设，再写为什么这个指标能支持或否定假设。若指标和假设没有关系，就不要把它放进报告。

边界是：分锁必须定义锁顺序，否则死锁风险上升。边界要写在正文里，因为工程设计通常不是“会不会”，而是“在什么条件下值得”。一个机制可能在大输入、高并发、失败恢复中必要，却在单线程小工具中完全多余。能说清这个取舍，才说明真正理解了机制。

\topic{单点分析：条件变量谓词}

围绕条件变量谓词，不要先记定义，先问它解决哪一个具体卡点：为什么 wait 必须放在 while 谓词里。在本章的有界队列、共享计数表和线程池状态里，这个卡点通常不会独立出现，而是和输入规模、数据分布、执行顺序、资源上限或故障注入一起出现。

朴素版本是：收到 notify 就认为条件成立。它在小输入下往往能通过，因为小输入没有把隐藏假设压穿。但是教材不能停在“它错了”这句话上，要说明它为什么一开始会被写出来：它减少了状态，靠调用栈维持顺序，靠当前机器的资源余量掩盖等待，靠人工观察代替指标。

一旦规模或失败形状变化，就会出现：虚假唤醒、竞争唤醒和关闭状态会让线程醒来后条件仍不成立。排查这类信号时，先不要改实现。先把它翻译成不变量问题：是不是同一个状态被多个拥有者写，是否某个输出位置没有唯一归属，是否某个对象生命周期跨过了异步边界，是否某个提交点缺少身份和校验。

机制模型是：条件变量等待的是状态谓词，不是通知本身。把模型落到代码时，至少要能指出一个字段、一个状态转换、一个测试和一个指标。字段让语义进入结构体，状态转换让路径可审查，测试让反例固定，指标让性能判断不靠猜。

实验观察是：构造多消费者和 close 场景验证。这里要避免一个常见错误：看到工具输出后直接写结论。正确做法是先写假设，再写为什么这个指标能支持或否定假设。若指标和假设没有关系，就不要把它放进报告。

边界是：notify 不是消息队列，不能携带业务事件。边界要写在正文里，因为工程设计通常不是“会不会”，而是“在什么条件下值得”。一个机制可能在大输入、高并发、失败恢复中必要，却在单线程小工具中完全多余。能说清这个取舍，才说明真正理解了机制。

\topic{单点分析：丢唤醒}

围绕丢唤醒，不要先记定义，先问它解决哪一个具体卡点：先 notify 后 wait 会不会丢事件。在本章的有界队列、共享计数表和线程池状态里，这个卡点通常不会独立出现，而是和输入规模、数据分布、执行顺序、资源上限或故障注入一起出现。

朴素版本是：把通知当成存储的信号。它在小输入下往往能通过，因为小输入没有把隐藏假设压穿。但是教材不能停在“它错了”这句话上，要说明它为什么一开始会被写出来：它减少了状态，靠调用栈维持顺序，靠当前机器的资源余量掩盖等待，靠人工观察代替指标。

一旦规模或失败形状变化，就会出现：若状态没有记录，等待者可能永远睡眠。排查这类信号时，先不要改实现。先把它翻译成不变量问题：是不是同一个状态被多个拥有者写，是否某个输出位置没有唯一归属，是否某个对象生命周期跨过了异步边界，是否某个提交点缺少身份和校验。

机制模型是：正确模式是先修改受锁保护状态，再 notify，让等待者检查状态。把模型落到代码时，至少要能指出一个字段、一个状态转换、一个测试和一个指标。字段让语义进入结构体，状态转换让路径可审查，测试让反例固定，指标让性能判断不靠猜。

实验观察是：用压力测试和超时测试暴露悬挂。这里要避免一个常见错误：看到工具输出后直接写结论。正确做法是先写假设，再写为什么这个指标能支持或否定假设。若指标和假设没有关系，就不要把它放进报告。

边界是：靠 sleep 修复丢唤醒是错误做法。边界要写在正文里，因为工程设计通常不是“会不会”，而是“在什么条件下值得”。一个机制可能在大输入、高并发、失败恢复中必要，却在单线程小工具中完全多余。能说清这个取舍，才说明真正理解了机制。

\topic{单点分析：信号量}

围绕信号量，不要先记定义，先问它解决哪一个具体卡点：资源槽位为什么不一定需要完整队列锁。在本章的有界队列、共享计数表和线程池状态里，这个卡点通常不会独立出现，而是和输入规模、数据分布、执行顺序、资源上限或故障注入一起出现。

朴素版本是：用 mutex 加计数器手写等待。它在小输入下往往能通过，因为小输入没有把隐藏假设压穿。但是教材不能停在“它错了”这句话上，要说明它为什么一开始会被写出来：它减少了状态，靠调用栈维持顺序，靠当前机器的资源余量掩盖等待，靠人工观察代替指标。

一旦规模或失败形状变化，就会出现：容易遗漏释放路径和异常安全。排查这类信号时，先不要改实现。先把它翻译成不变量问题：是不是同一个状态被多个拥有者写，是否某个输出位置没有唯一归属，是否某个对象生命周期跨过了异步边界，是否某个提交点缺少身份和校验。

机制模型是：信号量表示可用资源数量，适合连接数、缓冲槽和并发限流。把模型落到代码时，至少要能指出一个字段、一个状态转换、一个测试和一个指标。字段让语义进入结构体，状态转换让路径可审查，测试让反例固定，指标让性能判断不靠猜。

实验观察是：测试 acquire、release、超时和异常路径。这里要避免一个常见错误：看到工具输出后直接写结论。正确做法是先写假设，再写为什么这个指标能支持或否定假设。若指标和假设没有关系，就不要把它放进报告。

边界是：信号量不表达队列内容所有权，仍需数据结构同步。边界要写在正文里，因为工程设计通常不是“会不会”，而是“在什么条件下值得”。一个机制可能在大输入、高并发、失败恢复中必要，却在单线程小工具中完全多余。能说清这个取舍，才说明真正理解了机制。

\topic{单点分析：死锁}

围绕死锁，不要先记定义，先问它解决哪一个具体卡点：两个锁都正确为什么组合后会卡死。在本章的有界队列、共享计数表和线程池状态里，这个卡点通常不会独立出现，而是和输入规模、数据分布、执行顺序、资源上限或故障注入一起出现。

朴素版本是：每个函数按自己方便的顺序加锁。它在小输入下往往能通过，因为小输入没有把隐藏假设压穿。但是教材不能停在“它错了”这句话上，要说明它为什么一开始会被写出来：它减少了状态，靠调用栈维持顺序，靠当前机器的资源余量掩盖等待，靠人工观察代替指标。

一旦规模或失败形状变化，就会出现：线程之间形成等待环。排查这类信号时，先不要改实现。先把它翻译成不变量问题：是不是同一个状态被多个拥有者写，是否某个输出位置没有唯一归属，是否某个对象生命周期跨过了异步边界，是否某个提交点缺少身份和校验。

机制模型是：死锁来自资源获取顺序和不可抢占等待。把模型落到代码时，至少要能指出一个字段、一个状态转换、一个测试和一个指标。字段让语义进入结构体，状态转换让路径可审查，测试让反例固定，指标让性能判断不靠猜。

实验观察是：画等待图，定义全局锁顺序。这里要避免一个常见错误：看到工具输出后直接写结论。正确做法是先写假设，再写为什么这个指标能支持或否定假设。若指标和假设没有关系，就不要把它放进报告。

边界是：try lock 和超时只能缓解，不能替代设计。边界要写在正文里，因为工程设计通常不是“会不会”，而是“在什么条件下值得”。一个机制可能在大输入、高并发、失败恢复中必要，却在单线程小工具中完全多余。能说清这个取舍，才说明真正理解了机制。

\topic{单点分析：异常安全}

围绕异常安全，不要先记定义，先问它解决哪一个具体卡点：临界区里抛异常会不会留下半更新状态。在本章的有界队列、共享计数表和线程池状态里，这个卡点通常不会独立出现，而是和输入规模、数据分布、执行顺序、资源上限或故障注入一起出现。

朴素版本是：手动 lock 后在多个分支 unlock。它在小输入下往往能通过，因为小输入没有把隐藏假设压穿。但是教材不能停在“它错了”这句话上，要说明它为什么一开始会被写出来：它减少了状态，靠调用栈维持顺序，靠当前机器的资源余量掩盖等待，靠人工观察代替指标。

一旦规模或失败形状变化，就会出现：异常或早返回导致锁不释放或状态半更新。排查这类信号时，先不要改实现。先把它翻译成不变量问题：是不是同一个状态被多个拥有者写，是否某个输出位置没有唯一归属，是否某个对象生命周期跨过了异步边界，是否某个提交点缺少身份和校验。

机制模型是：RAII 和先构造后提交能保持异常路径清楚。把模型落到代码时，至少要能指出一个字段、一个状态转换、一个测试和一个指标。字段让语义进入结构体，状态转换让路径可审查，测试让反例固定，指标让性能判断不靠猜。

实验观察是：写抛异常的单元测试验证状态不变量。这里要避免一个常见错误：看到工具输出后直接写结论。正确做法是先写假设，再写为什么这个指标能支持或否定假设。若指标和假设没有关系，就不要把它放进报告。

边界是：异常安全不是只靠 lock guard，状态更新顺序也要设计。边界要写在正文里，因为工程设计通常不是“会不会”，而是“在什么条件下值得”。一个机制可能在大输入、高并发、失败恢复中必要，却在单线程小工具中完全多余。能说清这个取舍，才说明真正理解了机制。

\topic{单点分析：锁性能路径}

围绕锁性能路径，不要先记定义，先问它解决哪一个具体卡点：锁慢是因为内核调用还是 cache line 争用。在本章的有界队列、共享计数表和线程池状态里，这个卡点通常不会独立出现，而是和输入规模、数据分布、执行顺序、资源上限或故障注入一起出现。

朴素版本是：看到 mutex 就认为一定进内核。它在小输入下往往能通过，因为小输入没有把隐藏假设压穿。但是教材不能停在“它错了”这句话上，要说明它为什么一开始会被写出来：它减少了状态，靠调用栈维持顺序，靠当前机器的资源余量掩盖等待，靠人工观察代替指标。

一旦规模或失败形状变化，就会出现：低竞争 fast path 可能很快，高竞争才进入等待和唤醒。排查这类信号时，先不要改实现。先把它翻译成不变量问题：是不是同一个状态被多个拥有者写，是否某个输出位置没有唯一归属，是否某个对象生命周期跨过了异步边界，是否某个提交点缺少身份和校验。

机制模型是：锁性能由竞争、临界区长度、唤醒策略和 cache line 所有权共同决定。把模型落到代码时，至少要能指出一个字段、一个状态转换、一个测试和一个指标。字段让语义进入结构体，状态转换让路径可审查，测试让反例固定，指标让性能判断不靠猜。

实验观察是：记录等待次数、持锁时间和上下文切换。这里要避免一个常见错误：看到工具输出后直接写结论。正确做法是先写假设，再写为什么这个指标能支持或否定假设。若指标和假设没有关系，就不要把它放进报告。

边界是：不要为了避免 mutex 直接写错误 atomic 协议。边界要写在正文里，因为工程设计通常不是“会不会”，而是“在什么条件下值得”。一个机制可能在大输入、高并发、失败恢复中必要，却在单线程小工具中完全多余。能说清这个取舍，才说明真正理解了机制。

\topic{案例三轮推导：有界队列 close}

第一轮只写 reference。输入场景是：生产者和消费者都可能在关闭时等待。reference 的代码允许慢，甚至可以很笨，但语义必须清楚。它要说明哪些输入合法，哪些输入算错误，输出如何排序或比较，错误如何计数，重复运行是否得到同一摘要。

第二轮写朴素工程版本：只设置 closed 标志。这一版的价值是暴露直觉缺陷。写作时要诚实说明它为什么诱人：代码短，状态少，不需要额外结构，小样本容易通过。然后立刻给出反例，而不是直接跳到最终方案。

第三轮才写改进版本：在锁内更新状态并 notify all，让 push 和 pop 都检查谓词。改进不是装饰，而是针对第二轮反例的最小修复。如果修复引入了新的成本，也要同时写出来。例如局部合并会增加最终 merge，scan 会增加一轮 pass，manifest 会增加持久化开销，分片会增加元数据。

验收是：验证等待中关闭、空队列关闭、满队列关闭。验收报告还要额外写一条“不适用条件”。例如输入太小时，复杂机制可能不值得；数据分布变化时，热点处理可能失效；设备不支持某些计数器时，性能结论只能保留为假设。

\topic{案例三轮推导：分片锁 map}

第一轮只写 reference。输入场景是：多个 key 更新同一张计数表。reference 的代码允许慢，甚至可以很笨，但语义必须清楚。它要说明哪些输入合法，哪些输入算错误，输出如何排序或比较，错误如何计数，重复运行是否得到同一摘要。

第二轮写朴素工程版本：整张表一个 mutex。这一版的价值是暴露直觉缺陷。写作时要诚实说明它为什么诱人：代码短，状态少，不需要额外结构，小样本容易通过。然后立刻给出反例，而不是直接跳到最终方案。

第三轮才写改进版本：按 shard 加锁并定义合并阶段。改进不是装饰，而是针对第二轮反例的最小修复。如果修复引入了新的成本，也要同时写出来。例如局部合并会增加最终 merge，scan 会增加一轮 pass，manifest 会增加持久化开销，分片会增加元数据。

验收是：比较热点 key 和均匀 key 下的争用。验收报告还要额外写一条“不适用条件”。例如输入太小时，复杂机制可能不值得；数据分布变化时，热点处理可能失效；设备不支持某些计数器时，性能结论只能保留为假设。

\topic{案例三轮推导：双锁转移}

第一轮只写 reference。输入场景是：把任务从一个队列移动到另一个队列。reference 的代码允许慢，甚至可以很笨，但语义必须清楚。它要说明哪些输入合法，哪些输入算错误，输出如何排序或比较，错误如何计数，重复运行是否得到同一摘要。

第二轮写朴素工程版本：先锁源再锁目标或按调用者顺序锁。这一版的价值是暴露直觉缺陷。写作时要诚实说明它为什么诱人：代码短，状态少，不需要额外结构，小样本容易通过。然后立刻给出反例，而不是直接跳到最终方案。

第三轮才写改进版本：按地址或 id 定义全局锁顺序。改进不是装饰，而是针对第二轮反例的最小修复。如果修复引入了新的成本，也要同时写出来。例如局部合并会增加最终 merge，scan 会增加一轮 pass，manifest 会增加持久化开销，分片会增加元数据。

验收是：压力测试不应悬挂。验收报告还要额外写一条“不适用条件”。例如输入太小时，复杂机制可能不值得；数据分布变化时，热点处理可能失效；设备不支持某些计数器时，性能结论只能保留为假设。

\topic{实验手册：从命令到结论}

本章实验覆盖的机制包括：不变量、临界区粒度、条件变量谓词、丢唤醒、信号量、死锁、异常安全、锁性能路径。实验不是把这些名字逐个跑一遍，而是从同一条主线中抽取变量。先固定数据，再固定环境，再选择一个变量。变量可以是线程数、批大小、布局、分区函数、队列容量、故障注入点或重试预算，但一次只动一个。

实验输入建议分四类。正常输入用于确认功能没有被改坏；边界输入用于打到 off-by-one、空输入、满队列、零命中、重复消息等路径；压力输入用于暴露缓存、TLB、锁、队列、I/O 或网络成本；错误输入用于检查恢复、取消、幂等和错误分类。

每次实验至少记录五类结果。第一是 correctness digest，证明结果可信；第二是 throughput 或 elapsed time，说明端到端效果；第三是资源指标，例如内存峰值、队列水位、文件数量或 in-flight 数；第四是等待或失败指标，例如锁等待、超时、重试、late reply；第五是环境信息，例如 CPU、内核、编译器、输入规模和运行命令。

\begin{lstlisting}[numbers=none]
report skeleton:
  problem: one concrete symptom
  hypothesis: one boundary or cost source
  baseline: reference and naive version
  change: one mechanism only
  data: raw samples and digest
  result: what changed and what did not
  limit: unverified environment or counterexample
\end{lstlisting}

结论必须包含“没有变化”的部分。如果吞吐提高但内存峰值上升，要写；如果平均延迟降低但 p99 变差，要写；如果正确性通过但恢复路径未验证，要写。高质量教材不是让所有结果看起来漂亮，而是让读者学会在证据边界内说话。

\topic{Linux 源码映射：对象、状态和等待}

本章推荐源码入口是 \filepath{kernel/futex/core.c}、\filepath{kernel/locking/mutex.c}、\filepath{kernel/locking/semaphore.c}。阅读时不要展开成百科式笔记，而要围绕一个最小用户态现象。比如线程为什么睡眠，缺页为什么发生，写文件为什么没有立刻持久化，计数器为什么需要权限，队列为什么会唤醒。

源码阅读的核心是对象映射。用户态的线程，在内核里要找到 task 和调度状态；用户态的内存访问，要找到 VMA、页表项和 page；用户态的文件写入，要找到 file、inode、page cache 和 writeback；用户态的等待，要找到 wait queue、futex 或完成事件。

读内核代码时要特别注意状态字段如何变化。一个对象从 runnable 到 sleeping，从 clean page 到 dirty page，从 unlocked 到 contended，从 pending I/O 到 completed I/O，这些状态变化比函数名更能解释现象。把状态变化和本章实验指标对应起来，源码阅读才不是资料堆砌。

当前设备和源码版本可能不完全一致。报告要写明内核版本、阅读的源码版本、哪些结论来自源码机制、哪些结论来自本机实验。这能避免把源码阅读当作不可质疑的证明，也能训练读者区分机制理解和运行时证据。

\topic{贯穿项目验收：把能力写进系统}

本章进入贯穿项目时，不能只留下文字。至少要有一个目录、一个输入生成脚本或固定样本、一个 reference、一个朴素版本、一个改进版本、一个实验报告。代码可以小，但边界必须完整。

\begin{lstlisting}[numbers=none]
acceptance checklist:
  1. 为 bounded queue 写不变量注释
  2. 实现 close 语义和 notify all
  3. 加入锁等待指标
  4. 为分片计数表比较全局锁和分片锁
  5. 写死锁等待图练习
\end{lstlisting}

每个检查点都要对应一个失败路径。比如某个检查点是队列容量，就要能触发队列满；某个检查点是 manifest，就要能触发半写文件；某个检查点是 route epoch，就要能触发旧请求。没有失败注入的能力很容易只是正常路径演示。

验收时还要写回退策略。若改进版本复杂度太高、收益不足或设备环境不支持，项目应该能退回朴素但正确的版本。能回退说明边界清楚，不能回退往往说明机制已经和业务逻辑缠在一起。

本章结束前，读者应写一页复盘：本章新增了什么能力，解决了什么事故，引入了什么成本，下一章会复用哪个边界。这页复盘能让整本书保持连续，而不是每章都从零开始。

\topic{本节收束：回到最初事故}

最后再回到开头的事故：有界队列在关闭时偶发卡死，另一个计数表版本虽然加了锁，却在线程数增加后几乎不扩展。。如果现在重新排查，读者不应只说“可能是某机制的问题”，而应能写出一条可执行路线：固定输入，运行 reference，复现朴素失败，选择一个机制，改一个边界，采集指标，写出反例和结论限制。

这种路线就是本章真正要交付的能力。知识点会随着硬件、内核和框架变化而更新，但从事故出发、沿边界推导、用证据收束的习惯不会过时。

\topic{工程补充：实现走读与反例库}

把本章写成教材，最终要能落到一段能维护的工程实现。围绕有界队列、共享计数表和线程池状态，实现顺序建议从有界队列 close开始，因为它最容易把本章主失败暴露出来。先做 reference，再做朴素版本，最后做改进版本。reference 不追求速度，但它定义输入合同和输出摘要；朴素版本保留错误直觉；改进版本只改一个边界。

第一天实现不要碰太多机制。只写固定输入、固定输出和最小指标。输入要能触发这个事故：有界队列在关闭时偶发卡死，另一个计数表版本虽然加了锁，却在线程数增加后几乎不扩展。。输出摘要不要只用一个总数，要至少包含记录数、错误数、分片数、关键输出 checksum 和阶段状态。这样后面每次改布局、线程、队列或协议，都能回到同一个摘要比较。

第二天再引入分片锁 map。这一类案例通常负责暴露资源或调度问题。写实现时要特别注意观测点放在哪里：放在热循环里会改变成本，放在边界处又可能看不到细节。教学项目可以先用较粗粒度的阶段指标，等路径稳定后再增加更细的 trace。

第三天引入双锁转移。这一类案例通常负责恢复、尾部、长尾或提交边界。此时不要让改进版本直接覆盖朴素版本。两个版本应能在同一输入上并排运行，并输出同一套摘要。并排运行可以让读者看到差异来自机制，而不是来自输入变化或测量变化。

实现走读还要保持模块边界。输入生成、reference、朴素版本、改进版本、指标采集、报告输出最好分成独立函数或独立小文件。不要为了写得快，把实验、业务逻辑和日志打印都塞进一个函数。这样的代码短期能跑，长期无法解释，也无法被后续章节复用。

\topic{状态走读：哪些事实必须显式记录}

有界队列、共享计数表和线程池状态里最容易出错的地方，是把状态留在脑子里或调用栈里。只要出现线程、异步、重试、I/O 或分布式边界，调用栈就不再足够。教材要逼读者把事实写成字段：当前处理到哪里，谁拥有数据，哪次尝试有效，哪个输出已提交，哪个错误可重试，哪个指标能解释等待。

围绕不变量，要记录的是语义事实。它回答“结果为什么可信”。如果这个事实没有字段承载，测试只能在结果出错后才发现问题；如果字段存在，状态机就能在错误路径上提前拒绝非法转换。

围绕临界区粒度，要记录的是成本事实。它回答“为什么慢或为什么不扩展”。成本事实不一定进入业务结构，但必须进入实验报告。例如队列水位、批大小、任务等待、cache miss、重试次数、partition 大小，这些数字决定了优化方向。

围绕条件变量谓词，要记录的是边界事实。它回答“什么时候可以进入下一阶段”。边界事实尤其适合写成枚举状态，而不是一组布尔值。多个布尔值会产生 impossible state，枚举状态能让代码审查直接检查允许的转移。

状态走读的输出可以是一张简单表：状态名、拥有者、进入条件、退出条件、持久化要求、观测指标。读者不需要一开始画出完美状态机，但必须知道哪些状态不能靠注释维持。凡是会跨线程、跨文件、跨消息或跨重启的状态，都应优先显式化。

\begin{lstlisting}[numbers=none]
state review table:
  state name
  owner
  allowed incoming events
  allowed outgoing events
  persisted or transient
  metric that proves progress
\end{lstlisting}

\topic{反例库：防止机制变成口号}

反例库的作用，是防止读者把本章机制学成一句正确但空泛的话。每个反例都要小到可以复现，大到能代表真实系统的一类失败。反例不一定导致崩溃，很多时候它只表现为吞吐不上升、输出顺序不稳定、恢复后重复、队列增长或长尾放大。

反例：如果只按不变量的直觉写代码，而不检查“锁到底保护变量还是保护关系”，就会落回朴素路径“哪个变量会被多线程访问就给哪个变量加锁”。触发条件可以设计成：扩大输入、改变分布、插入等待、打乱顺序、制造重复或提前关闭。预期失败信号是：多个变量之间的关系可能在锁外被破坏。这个反例应成为测试或实验的一部分，而不是只写在正文里。

反例：如果只按临界区粒度的直觉写代码，而不检查“一个大锁简单，为什么有时不够好”，就会落回朴素路径“用一个 mutex 保护整个对象”。触发条件可以设计成：扩大输入、改变分布、插入等待、打乱顺序、制造重复或提前关闭。预期失败信号是：低竞争时清楚，高竞争时所有操作串行。这个反例应成为测试或实验的一部分，而不是只写在正文里。

反例：如果只按条件变量谓词的直觉写代码，而不检查“为什么 wait 必须放在 while 谓词里”，就会落回朴素路径“收到 notify 就认为条件成立”。触发条件可以设计成：扩大输入、改变分布、插入等待、打乱顺序、制造重复或提前关闭。预期失败信号是：虚假唤醒、竞争唤醒和关闭状态会让线程醒来后条件仍不成立。这个反例应成为测试或实验的一部分，而不是只写在正文里。

反例：如果只按丢唤醒的直觉写代码，而不检查“先 notify 后 wait 会不会丢事件”，就会落回朴素路径“把通知当成存储的信号”。触发条件可以设计成：扩大输入、改变分布、插入等待、打乱顺序、制造重复或提前关闭。预期失败信号是：若状态没有记录，等待者可能永远睡眠。这个反例应成为测试或实验的一部分，而不是只写在正文里。

反例：如果只按信号量的直觉写代码，而不检查“资源槽位为什么不一定需要完整队列锁”，就会落回朴素路径“用 mutex 加计数器手写等待”。触发条件可以设计成：扩大输入、改变分布、插入等待、打乱顺序、制造重复或提前关闭。预期失败信号是：容易遗漏释放路径和异常安全。这个反例应成为测试或实验的一部分，而不是只写在正文里。

反例：如果只按死锁的直觉写代码，而不检查“两个锁都正确为什么组合后会卡死”，就会落回朴素路径“每个函数按自己方便的顺序加锁”。触发条件可以设计成：扩大输入、改变分布、插入等待、打乱顺序、制造重复或提前关闭。预期失败信号是：线程之间形成等待环。这个反例应成为测试或实验的一部分，而不是只写在正文里。

反例库还要写出修复后如何证明不再发生。证明不能只靠“这次没有复现”。更好的方式是把反例输入固定下来，把状态摘要固定下来，把失败路径上的关键指标固定下来。以后重构时，这些反例就是回归测试。

\topic{实验协议：同一脚本跑出四类结论}

本章实验脚本要围绕一句话展开：容量设为一，同时启动多个生产者和消费者；在空、满、半满、等待中分别调用 close。。脚本至少要支持四种模式：correctness、pressure、fault、report。correctness 模式跑小输入和边界输入，pressure 模式跑大输入或高并发，fault 模式启用故障注入，report 模式把原始样本整理成表。

correctness 模式不应该被性能优化绕过。它要在每个版本后运行，覆盖空输入、单元素、非整齐长度、重复 key、非法记录、关闭边界和恢复边界。如果某章没有这些具体路径，也要选择与本章机制等价的边界输入。

pressure 模式要把瓶颈逼出来。输入太小，很多机制看起来都没有意义；输入过大，又可能被外部环境噪声掩盖。合适的做法是用几档规模扫描，找到性能曲线的拐点，再围绕拐点解释。只报告最大规模，往往看不到瓶颈迁移过程。

fault 模式要把错误变成输入。错误不应依赖人工按 Ctrl-C 或碰运气，而应在脚本里用固定种子、固定时刻、固定任务或固定文件触发。例如在写 manifest 前杀进程，在某个消息上丢响应，在某个 partition 上制造热点，在某个 worker 上注入慢 I/O。

report 模式要保留原始样本，而不是只保留平均值。报告可以生成 CSV、纯文本摘要或 LaTeX 表格，但至少要能重新计算平均、分位数、最大值和错误计数。如果读者只能看到最终一句结论，就无法判断结论是否可靠。

\begin{lstlisting}[numbers=none]
lab modes:
  correctness: small and boundary inputs
  pressure: scale, concurrency, and distribution sweep
  fault: deterministic failure injection
  report: raw samples, digest, and conclusion limits
\end{lstlisting}

\topic{源码审查：从实现反推本章原则}

源码审查时不要先看函数数量，也不要先看是否用了某个高级技巧。先从数据入口开始，一路跟到提交或输出。每经过一个边界，就问这个边界有没有字段、状态、测试和指标。这个方法比逐行读代码更有效，因为它直接对应系统失败路径。

以本章为例，最小走查路线是：遇到卡死先画等待图，标出每个线程持有什么锁、等待什么谓词；不要用 sleep 或超时掩盖丢唤醒。。这条路线里，每个动作都应该能在源码中找到对应位置。如果找不到，说明代码把语义藏在约定里；如果能找到但没有测试，说明边界没有被固定；如果有测试但没有指标，说明性能和恢复结论仍然靠猜。

审查还要看错误路径。很多项目的正常路径清晰，错误路径却散在多个 return、异常、回调和清理函数里。教材里的示例必须展示错误如何返回、资源如何释放、等待者如何唤醒、临时输出如何清理、重复请求如何去重。否则读者学到的只是 happy path。

最后看命名。状态字段、计数器、日志事件和文件名都应体现语义身份。一个叫 data、flag、tmp、done 的字段，在小程序里能理解，在系统里会变成歧义源。更好的名字应包含对象、阶段和语义，例如 \code{current_attempt}、\code{committed_checksum}、\code{producer_wait_ns}、\code{route_epoch}。

\topic{迁移说明：从本章到下一章}

本章内容不能学完就断掉。每个机制都要问它会在后续哪里复用。单核里的依赖链会迁移到 SIMD 和并行归约；cache line 所有权会迁移到锁、原子和分片计数；有界队列会迁移到 I/O 背压和分布式流控；manifest 会迁移到 shuffle 和 checkpoint。

对于有界队列、共享计数表和线程池状态，可迁移的是语义边界和实验方法，不可直接迁移的是具体数字。吞吐、延迟、批大小、线程数、队列容量、partition 数量，都要在新输入和新机器上重新测。这也是为什么本书反复要求记录环境和原始样本。

迁移时还要防止过度抽象。看到两个章节都有“队列”，不代表它们可以共用同一个类；看到两个章节都有“状态机”，不代表可以写一个万能状态机框架。抽象只有在减少真实重复、保护共同不变量、并且不会隐藏关键成本时才值得引入。

读者可以在每章末尾写三句话。第一，本章最重要的不变量是什么。第二，本章最可能迁移到后续哪一章。第三，本章最不能照搬的性能数字是什么。这三句话能帮助整本书形成连续推理，而不是把章节当作孤立文章。

\topic{章末审查：验收题、反例输入和提交标准}

本章最后再做一次审查。审查不是重复正文，而是把有界队列、共享计数表和线程池状态压成可检查的交付物。最小验收题应覆盖这些案例：有界队列 close、分片锁 map、双锁转移。如果读者只能描述概念，却不能为这些案例写出 reference、朴素版本、反例和改进版本，就还没有达到本章要求。

第一类验收题检查语义。给定一组固定输入，要求写出串行 reference 和结果摘要。摘要要能暴露遇到卡死先画等待图，标出每个线程持有什么锁、等待什么谓词；不要用 sleep 或超时掩盖丢唤醒。。答案必须说明哪些字段参与比较，哪些错误要计数，哪些状态可以忽略，哪些状态必须持久化或记录。

第二类验收题检查机制推导。要求从不变量、临界区粒度、条件变量谓词、丢唤醒、信号量中选择一个机制，先写朴素方案为什么自然，再写它在本章主线中如何失败，最后写一个只改变单一边界的改进。答案不能直接给最终方案。没有朴素版本，就没有推导；没有失败证据，就没有必要性；没有反例，就没有边界。

第三类验收题检查实验设计。要求给出正常输入、边界输入、压力输入和错误输入四组样本。每组样本都要说明目的：验证功能、打到边界、暴露成本、检查恢复。如果一个样本只能让程序跑完，却不能支持或否定某个假设，就不应进入实验矩阵。

\begin{lstlisting}[numbers=none]
chapter review checklist:
  reference digest is stable
  naive version is preserved
  counterexample input is committed
  improved version changes one boundary
  metrics explain the chosen hypothesis
  report states unverified limits
\end{lstlisting}

\topic{反例输入怎么设计}

反例输入不要追求稀奇，而要追求直击本章边界。边界输入通常比大输入更有价值，因为它能把 off-by-one、空队列、重复请求、旧 attempt、半写文件、错误路由、尾部元素、分片倾斜这些问题直接打出来。

围绕有界队列、共享计数表和线程池状态，一个好的反例输入应满足三个条件。第一，人工可以解释预期结果；第二，朴素版本会稳定暴露问题；第三，改进版本通过后仍能说明机制边界。如果反例太随机，失败了也难解释；如果反例太简单，可能无法触发真实风险。

反例还要保留原始材料。输入文件、随机种子、故障注入配置、运行命令、环境摘要和输出 digest 都应进入仓库或报告。不要只在正文里描述“构造一个错误输入”。没有材料，下一次重构时无法回归。

\topic{项目提交前的自检}

提交本章项目代码前，先问六个问题。第一，reference 是否还在。第二，朴素版本是否还能跑。第三，改进版本是否只改变一个主要机制。第四，错误路径是否有测试。第五，指标是否能解释结论。第六，报告是否写明没有验证的硬件或系统边界。

如果这些问题有任何一个回答不清，就先不要把章节标为完成。系统书的质量来自可复查，而不是来自一次性写出大量文字。读者能复查，作者能复查，后续章节也能复用，才说明本章真正稳定。

最后回到本章事故：有界队列在关闭时偶发卡死，另一个计数表版本虽然加了锁，却在线程数增加后几乎不扩展。。章末自检的意义，就是让读者面对同类事故时能不慌乱。先固定事实，再找第一个分歧，再选机制，再写实验。这个顺序会在整本第二册反复出现。

\

t

o

p

i

c

{

学

习

复

盘

：

把

本

章

能

力

带

到

下

一

章

}







离

开

本

章

前

，

读

者

应

能

把

有

界

队

列

、

共

享

计

数

表

和

线

程

池

状

态

讲

成

一

条

完

整

因

果

链

，

而

不

是

讲

成

一

组

概

念

。

起

点

是

事

故

：

有

界

队

列

在

关

闭

时

偶

发

卡

死

，

另

一

个

计

数

表

版

本

虽

然

加

了

锁

，

却

在

线

程

数

增

加

后

几

乎

不

扩

展

。

；

中

间

是

朴

素

方

案

、

失

败

路

径

、

状

态

字

段

和

实

验

矩

阵

；

终

点

是

能

进

入

贯

穿

项

目

的

代

码

边

界

。

如

果

复

述

时

只

能

说

出

术

语

，

却

说

不

清

第

一

个

分

歧

出

现

在

输

入

、

执

行

、

同

步

、

I

/

O

、

提

交

还

是

恢

复

，

就

应

该

回

到

本

章

案

例

重

新

走

一

遍

。







本

章

最

少

要

带

走

两

个

能

力

。

第

一

个

是

围

绕

不

变

量

建

立

语

义

边

界

：

哪

些

输

入

合

法

，

哪

些

状

态

可

见

，

哪

些

结

果

可

以

比

较

，

哪

些

错

误

必

须

外

显

。

第

二

个

是

围

绕

临

界

区

粒

度

建

立

成

本

边

界

：

哪

些

字

节

被

搬

动

，

哪

些

线

程

在

等

待

，

哪

些

队

列

在

积

压

，

哪

些

重

试

或

失

败

放

大

了

工

作

量

。

语

义

边

界

让

结

果

可

信

，

成

本

边

界

让

优

化

有

方

向

，

两

者

缺

一

不

可

。







本

章

案

例

包

括

有

界

队

列



c

l

o

s

e

、

分

片

锁



m

a

p

、

双

锁

转

移

。

复

盘

时

不

要

只

看

最

终

改

进

版

本

，

还

要

回

看

朴

素

版

本

为

什

么

会

被

写

出

来

。

很

多

工

程

事

故

的

根

源

不

是

作

者

不

知

道

高

级

机

制

，

而

是

小

规

模

条

件

成

立

后

，

没

有

在

规

模

、

并

发

、

故

障

和

恢

复

条

件

变

化

时

重

新

审

查

边

界

。

教

材

反

复

保

留

朴

素

版

本

，

就

是

为

了

训

练

这

种

迁

移

审

查

能

力

。







把

本

章

带

到

下

一

章

时

，

只

迁

移

方

法

，

不

照

搬

数

字

。

r

e

f

e

r

e

n

c

e

、

反

例

、

单

变

量

实

验

、

状

态

表

、

指

标

和

结

论

边

界

可

以

迁

移

；

具

体

吞

吐

、

延

迟

、

线

程

数

、

批

大

小

、

队

列

容

量

、

分

片

数

量

都

必

须

重

新

测

。

这

条

原

则

能

避

免

读

者

把

某

台

机

器

上

的

经

验

写

成

普

遍

规

律

。







最

后

给

自

己

留

一

个

三

句

话

笔

记

：

本

章

保

护

的

核

心

不

变

量

是

什

么

；

本

章

新

增

的

成

本

或

复

杂

度

是

什

么

；

本

章

哪

个

实

验

最

值

得

在

后

续

章

节

复

用

。

这

三

句

话

比

背

诵

术

语

更

重

要

，

因

为

它

能

把

章

节

串

成

一

套

工

程

判

断

。

\topic{本章质量标准}

读完本章，读者应该能围绕有界队列、共享计数表和线程池状态讲清完整推导：朴素方案为什么成立，哪条失败路径先出现，引入的机制保护了哪个不变量，实验如何证明或否定假设。

如果正文只列概念，没有具体输入、状态、时间线和反例，就不合格；如果只给最终代码，没有朴素版本和失败证据，也不合格；如果只给 benchmark 数字，没有语义 reference 和结论边界，同样不合格。

本章最终留下的是判断力：面对一个慢、错、卡死、重复提交或无法恢复的系统，先固定问题，再画边界，再写 reference，再做对照，再改机制。这个顺序比任何单个技巧更重要。

\section{本章落到贯穿项目}

Compute Systems Engine 当前的 \code{BoundedMpmcQueue} 已经有容量、head、tail、size、mutex 和 \code{not_full}，但它还不是完整线程池队列，因为它只有阻塞 push 和非阻塞 try pop，没有 \code{not_empty} 等待，也没有 close。作为教学阶段的最小模块可以接受，但后续运行时章节必须补全关闭语义和消费者等待，否则 worker 无法优雅阻塞和退出。

本章要求项目至少增加设计文档或测试计划：列出队列不变量，列出 push、try pop、未来 blocking pop 和 close 的状态转换，列出需要的测试。测试不能只 push 两个值再 pop 两个值，还要覆盖容量为一、队列满时生产者等待、队列空时消费者等待、close 唤醒等待者、多生产者多消费者压力、异常容量和析构时没有等待线程。

\topic{状态机比代码片段更重要}

同步代码最怕只看局部片段。一个 \code{push} 函数看起来正确，一个 \code{pop} 函数看起来正确，合在一起仍可能在关闭、异常、等待和析构时出错。因此并发容器应先写状态机。以有界队列为例，状态至少包括 open empty、open partial、open full、closing nonempty、closed empty。每个状态下，\code{push}、\code{try_push}、\code{pop}、\code{try_pop}、\code{close}、\code{wait_empty} 的行为都要定义。

状态机能暴露隐藏问题。open full 状态下，阻塞 \code{push} 等待什么谓词？如果此时调用 \code{close}，生产者应返回失败还是继续等待？closing nonempty 状态下，消费者能否继续取出剩余任务？closed empty 状态下，\code{pop} 返回空值还是抛异常？这些问题不回答，代码只能靠偶然行为。一个高质量并发库的接口文档，应像协议一样说明状态转换。

状态机还要说明线性化点。对于 \code{push}，元素写入槽位并更新 \code{tail} 和 \code{size} 后，操作对消费者可见；对于 \code{pop}，元素被取出并更新 \code{head} 和 \code{size} 后，槽位重新可用；对于 \code{close}，\code{closed_} 从 false 变为 true 的那一刻，后续生产者必须观察到关闭。线性化点不清楚，就无法判断并发调用的先后语义。

\topic{背压不是错误路径}

有界队列满了，生产者等待，这不是程序变差，而是背压在工作。背压的目的，是让下游资源不足时，上游不要无限制地产生任务和占用内存。没有背压的系统，短时间 benchmark 可能看起来更快，因为它把工作全部塞进内存；一旦输入持续增长，就会出现内存膨胀、延迟暴涨和最终崩溃。有界队列用等待或拒绝把压力反馈给上游。

背压策略要和业务语义匹配。日志系统可以选择丢弃低优先级调试日志；任务运行时通常不能丢计算任务；网络服务可以对新请求返回繁忙；分布式 shuffle 可以暂停上游分片读取。队列只提供机制，业务决定策略。教材中如果只写“队列满则阻塞”，还不够完整；还要讨论阻塞会不会占住 compute worker，是否需要异步等待，是否需要超时，是否需要取消。

背压指标也要写入报告。队列满的次数、生产者等待总时长、最大队列长度、丢弃任务数、拒绝请求数，这些指标能说明系统是在健康限流，还是下游长期跟不上。没有指标，开发者可能把背压误认为锁慢，然后盲目增加容量。容量变大只会把压力藏得更深，尾延迟和内存风险更高。

\topic{锁和异常传播}

线程池中的任务可能抛异常。异常传播如果处理不好，会破坏锁保护的不变量。最基本原则是：任务执行不应发生在队列锁内，任务异常不应让未完成计数漏减，异常对象的发布要有同步关系。队列只负责交付任务，worker 从队列取出任务后释放锁，再执行任务。执行结束后，不管成功还是异常，都要更新任务状态和未完成计数，并通知等待者。

C++ 的 RAII 能帮助释放锁，但不能自动维护业务状态。若一个任务执行前增加了 in-flight 计数，异常路径必须减少；若 worker 把错误写入 shared state，等待者读取错误前要有同步；若异常对象包含动态分配资源，生命周期要覆盖读取者。并发异常不是“catch 一下”这么简单，它和状态机、通知和对象生命周期同时相关。

一个实用模式是把任务结果分成状态和载荷。状态可以是 Pending、Running、Succeeded、Failed、Cancelled；载荷是返回值或错误信息。worker 先写载荷，再在锁内或 release 语义下发布状态，然后通知等待者。等待者醒来后检查状态，再读取载荷。这个模式和原子章节的对象发布一致，只是这里多了一层任务语义。

\topic{测试等待路径要制造极端条件}

同步代码的测试不能只跑正常路径。容量为一的队列比容量很大的队列更容易暴露满和空之间的状态转换；一个 producer 一个 consumer 能测基本交替；多个 producer 一个慢 consumer 能测 \code{not_full}；一个 producer 多个 consumer 能测 \code{not_empty} 和惊群；关闭时有线程阻塞在 push 和 pop，能测 \code{notify_all} 是否完整。小容量、慢消费者、随机 yield、重复关闭和异常任务，是并发测试的基本工具。

测试还应覆盖析构边界。线程池析构时是否自动 shutdown？如果还有外部线程提交任务，行为是什么？析构能否在 worker 正在执行任务时等待？如果析构发生在某个任务持有业务锁时，会不会死锁？这些问题不一定都由队列解决，但队列和线程池接口必须定义。未定义行为可以存在，但必须明确禁止，而不是让用户猜。

压力测试不能替代不变量证明，但能发现实现错误。可以把队列操作与串行 reference 对比：多个线程随机 push 和 pop，同时记录成功操作日志，最终验证元素没有丢失、没有重复、容量不越界、关闭后没有新元素进入。对于不可确定顺序的并发队列，验证重点不是全局顺序完全一致，而是每个成功 push 最多被 pop 一次，关闭语义符合文档。

\section{本章习题}

\begin{exercisebox}
习题一：为当前项目的 \code{BoundedMpmcQueue} 写不变量清单。逐行解释 \code{push} 和 \code{try_pop} 怎样保持这些不变量。指出当前接口缺少哪些线程池需要的语义。

习题二：把 bounded queue 设计成支持 \code{close}。不要求立刻写完整代码，但必须画出状态转换：open nonempty、open empty、open full、closed nonempty、closed empty。说明每个状态下 push 和 pop 应返回什么。

习题三：设计一个队列等待实验。多个生产者写入，消费者故意变慢，让队列经常满。记录生产者等待次数、最大队列长度和总吞吐。解释这是背压，不是单纯性能退化。

习题四：阅读 glibc、musl 或 Linux futex 文档和源码入口，写清 mutex 低竞争 fast path 和竞争 slow path 的差异。报告中必须说明哪些部分发生在用户态，哪些部分进入内核。
\end{exercisebox}

\begin{warningbox}
不要把锁竞争问题简单理解成“换成无锁”。如果临界区保护的是复杂不变量，无锁实现可能更难证明正确，也可能在高竞争下 CAS 失败更多。
\end{warningbox}

\begin{labbox}
实验：实现一个 bounded queue。先用 mutex 和 condition variable 写阻塞 push/pop，再记录高生产者数量下的等待时间。解释锁保护的状态和等待谓词。
\end{labbox}
